%% 
%% Author's note:
%% This was a template package which contains a lot of comments and code that were not written by me.
%%


%% Follow comments to support use.

%%%%%%%%%%%%%%%%%%%%%%%%%%%%%%%%%%%%%%%%%%%%%%%%%%%%%%%%%
%% STEP 1: Choose options for MSc / BSc / seminar layout and your bibliographic style
%%%%%%%%%%%%%%%%%%%%%%%%%%%%%%%%%%%%%%%%%%%%%%%%%%%%%%%%%

%%  Language: 
%%      finnish, swedish, or english
%%  Pagination (use twoside by default)  
%%      oneside or twoside,
%%  Study programme / kind of report
%%      csm  = Master's thesis in Computer Science Master's Programme;
%%      tkt = Bachelor's thesis in Computer Science Bachelor's Programme;
%%      seminar = seminar report
%%  For Master's thesis choose your line or track:
%%      (30 cr thesis, 2020 onwards, Master's Programme in Computer Science = csm)
%%      software-track-2020 = Software study track
%%      algorithms-track-2020 = Algorithms study track
%%      networking-track-2020 = Networking study track

\documentclass[english,twoside,censored,tkt,null]{HYthesisML} 


%% TODO: Check which is better (openright or openany).
% If wanted, open new chapters only at right page.
% By default, "openany".
%\PassOptionsToClass{openright,twoside,a4paper}{report}
\PassOptionsToClass{openany,twoside,a4paper}{report}

\usepackage{csquotes}
%%%%%%%%%%%%%%%%%%%%%%%%%%%%%%%%%%%%%%%%%%%%%%%%%%%%%%%%%
%% REFERENCES
%% Some notes on bibliography usage and options:
%% natbib -> you can use, e.g., \citep{} or \parencite{} for (Einstein, 1905); with APA \cite -> Einstein, 1905 without ()
%% maxcitenames=2 -> only 2 author names in text citations, if more -> et al. is used
%% maxbibnames=99 as no great need to suppress the biliography list in a thesis
%% for more information see biblatex package documentation, e.g., from https://ctan.org/pkg/biblatex 

%% Reference style: select one 
%% for APA = Harvard style = authoryear -> (Einstein, 1905) use:
\usepackage[style=authoryear,bibstyle=authoryear,backend=biber,natbib=true,maxnames=99,maxcitenames=2,uniquelist=minyear,giveninits=true,uniquename=mininit]{biblatex}
%% for numeric = Vancouver style -> [1] use:
%\usepackage[style=numeric,bibstyle=numeric,backend=biber,natbib=true,maxbibnames=99,giveninits=true,uniquename=init]{biblatex}
%% for alpahbetic -> [Ein05] use:
%\usepackage[style=alphabetic,bibstyle=alphabetic,backend=biber,natbib=true,maxbibnames=99,giveninits=true,uniquename=init]{biblatex}
%

\addbibresource{bibliography.bib}
% in case you want the final delimiter between authors & -> (Einstein & Zweistein, 1905) 
% \renewcommand{\finalnamedelim}{ \& }
% List the authors in the Bibilipgraphy as Lastname F, Familyname G,
\DeclareNameAlias{sortname}{family-given}
% remove the punctuation between author names in Bibliography 
%\renewcommand{\revsdnamepunct}{ }

% Asked chat.openai.com: I'm using biblatex package, but it prints odd format volume.number in reference list, while this should be volume (number). How to fix this?
%Answer (shortened):
% Redefine the volume and number format
%\DeclareFieldFormat[article]{volume}{\mkbibbold{#1}}
%\DeclareFieldFormat[article]{number}{\mkbibparens{#1}}
% Please remove the boldface from the volume and number fields:
\DeclareFieldFormat[article]{volume}{#1}
\DeclareFieldFormat[article]{number}{\mkbibparens{#1}}
% Now this generated format volume.(number). How to remove the dot? (and so on with some adjustments):
\renewbibmacro*{volume+number+eid}{%
  \setunit{\addcomma\space}% 
  \printfield{volume}%
  %\setunit*{\addnbspace\addthinspace}% <-- Adjust spacing here
  \printfield{number}%
  \setunit{\addcomma\space}%
  \printfield{eid}}


%% Block of definitions for fonts and packages for picture management.
%% In some systems, the figure packages may not be happy together.
%% Choose the ones you need.

%\usepackage[utf8]{inputenc} % For UTF8 support, in some systems. Use UTF8 when saving your file.

\usepackage{lmodern}         % Font package, again in some systems.
\usepackage{textcomp}        % Package for special symbols
\usepackage[pdftex]{color, graphicx} % For pdf output and jpg/png graphics
\usepackage{epsfig}
\usepackage{subfigure}
\usepackage[pdftex, plainpages=false]{hyperref} % For hyperlinks and pdf metadata
\usepackage{fancyhdr}        % For nicer page headers
\usepackage{tikz}            % For making vector graphics (hard to learn but powerful)
%\usepackage{wrapfig}        % For nice text-wrapping figures (use at own discretion)
\usepackage{amsmath, amssymb} % For better math

\singlespacing               %line spacing options; normally use single

\fussy
%\sloppy                      % sloppy and fussy commands can be used to avoid overlong text lines
% if you want to see which lines are too long or have too little stuff, comment out the following lines
% \overfullrule=1mm
% to see more info in the detailed log about under/overfull boxes...
% \showboxbreadth=50 
% \showboxdepth=50



%%%%%%%%%%%%%%%%%%%%%%%%%%%%%%%%%%%%%%%%%%%%%%%%%%%%%%%%%
%% STEP 2:
%%%%%%%%%%%%%%%%%%%%%%%%%%%%%%%%%%%%%%%%%%%%%%%%%%%%%%%%%
%% Set up personal information for the title page and the abstract form.
%% Replace parameters with your information.
\title{AI-Powered Social Engineering}


\author{Riku Talvisto}
\date{\today}

% Set supervisors, use the titles according to the thesis language
% in English Prof. or Dr., or in Finnish toht. or tri or FT, TkT, Ph.D. or in Swedish... 
\supervisors{Dr. Lea Kutvonen}

\keywords{social engineering, artificial intelligence, AI, cybersecurity, security, hacking, psychology}
\additionalinformation{\translate{\track}}

%% For seminar reports:
%%\additionalinformation{Name of the seminar}

%% Provide classification terms, to appear on the abstract page.
%% Replace the classification terms below with the ones that match your work.
%% ACM Digital library provides a taxonomy and a tool for classification
%% in computer science. Use 1-3 paths, and use right arrows between the
%% about three levels in the path; each path requires a new line.

%% TODO: Which are most suitable? Should I include one, 2 or 3? Should I include the document type (surveys and overviews) and is that correct?
%% TODO: Should any title be in bold or any in italics?
\classification{\protect{\ \\
\  Social and professional topics $\rightarrow$ Computing / technology policy  $\rightarrow$ Computer crime \\ $\rightarrow$ \textbf{Social engineering attacks} \\
%%\  Human-centered computing $\rightarrow$ Collaborative and social computing  $\rightarrow$ Collaborative and social computing theory, concepts and paradigms $\rightarrow$ Social engineering (social sciences) \\
\  Security and privacy $\rightarrow$ Intrusion/anomaly detection and malware mitigation \\ $\rightarrow$ \textbf{Social engineering attacks}
}}

%% If you want to quote someone special. You can comment this line out and there will be nothing on the document.
%\quoting{Bachelor's degrees make pretty good placemats if you get them laminated.}{Jeph Jacques}


%% TODO: Remember to make the final PDF a PDF/A compliant version for archival purposes!

%% OPTIONAL STEP: Set up properties and metadata for the pdf file that pdfLaTeX makes.
%% Your name, work title, and keywords are recommended.
\hypersetup{
    unicode=true,           % to show non-Latin characters in Acrobat’s bookmarks
    pdftoolbar=true,        % show Acrobat’s toolbar?
    pdfmenubar=true,        % show Acrobat’s menu?
    pdffitwindow=false,     % window fit to page when opened
    pdfstartview={FitH},    % fits the width of the page to the window
    pdftitle={},            % title
    pdfauthor={},           % author
    pdfsubject={},          % subject of the document
    pdfcreator={},          % creator of the document
    pdfproducer={pdfLaTeX}, % producer of the document
    pdfkeywords={something} {something else}, % list of keywords for
    pdfnewwindow=true,      % links in new window
    colorlinks=true,        % false: boxed links; true: colored links
    linkcolor=black,        % color of internal links
    citecolor=black,        % color of links to bibliography
    filecolor=magenta,      % color of file links
    urlcolor=cyan           % color of external links
}

%%-----------------------------------------------------------------------------------

\begin{document}

% Generate title page.
\maketitle

%%%%%%%%%%%%%%%%%%%%%%%%%%%%%%%%%%%%%%%%%%%%%%%%%%%%%%%%%
%% STEP 3:
%%%%%%%%%%%%%%%%%%%%%%%%%%%%%%%%%%%%%%%%%%%%%%%%%%%%%%%%%
%% Write your abstract in the separate file, to be positioned here.
%% You can make several abstract pages (if you want it in different languages),
%% in which case you should also define the language of the abstract,
%% as below.



    %%%%%%%%%%%%%%%%%%%%%%%%%%%%%%%%%%%%%%%%%%%%%%%
    %% ABSTRACT (English)                        %%
    %%%%%%%%%%%%%%%%%%%%%%%%%%%%%%%%%%%%%%%%%%%%%%%


\begin{otherlanguage}{english}
\begin{abstract}

\begin{comment}

Guides:
    - Max about 100-200 words
        - Some other theses have almost fully covered the available area
    - If the abstract page is too full, is that demotivating to the reader?
    - Study material states that the abstract text is short, usually just one paragraph
    - What has been studied
    - How it has been studied
    - What results have been observed

TODO:
    [ ] What is the research question
    [ ] What is social engineering, quick definition?
    [ ] How does modern AI augment SE attacks and countermeasures
    [ ] What is analyzed in this thesis
    [ ] Properly motivate the reader to read my thesis
    [ ] Explain that this is written in an easy-to-read manner?
    [ ] What results have been observed

What to cover:
    - What is SE?
    - Modern AI
    - Attacks
        - Deepfake synthetic content, videos, live voice morphing
        - Highly personalized phishing content (natural language processing)
        - Automated OSINT gathering
    - AI augments both attacks and countermeasures
    - Countermeasures
        - User awareness & training programs
        - Company policy & company culture
        - Real-time threat detection
        - Vulnerability detection
    - Emerging challenges

\end{comment}

This thesis explores the evolving landscape of social engineering (SE) in the age of modern artificial intelligence (AI). While AI offers modern opportunities for enhancing cybersecurity measures, it simultaneously empowers malicious actors with sophisticated tools for crafting highly targeted and effective social engineering attacks.

Delving into the various ways AI is being exploited to augment social engineering tactics, including the creation of highly believable synthetic media like deepfake images, videos and real-time voice morphing, the generation of personalized phishing messages through natural language processing.

Conversely, the thesis also explores how AI can be harnessed to bolster countermeasures by enabling real-time threat detection, identifying potential vulnerabilities and facilitating comprehensive employee training programs. By analyzing the dual-faceted impact of AI on SE, this thesis aims to provide a comprehensive understanding of the emerging challenges and opportunities in this domain.

Given the rapidly evolving nature of AI and its expanding capabilities, much of the content presented is speculative.These projections are grounded in analysis of established social engineering attacks and observed trends in AI development.

%This thesis explores the intersection of artificial intelligence (AI) and social engineering, examining the emerging threats as well as opportunities for defense in this space.

%The findings reveal that AI-powered attacks are more convincing and successful than traditional attacks, but AI-driven defense mechanisms can improve detection rates, and AI can develop more engaging training material(?).

% The integration of artificial intelligence (AI) into social engineering practices presents exceptional threats to privacy and security, necessitating a truly comprehensive re-evaluation, and enhancement, of current cybersecurity measures.

%The study contributtes to our understanding of AI's impact on social engineering and highlights the need for cybersecurity professionals to adapt to these emerging threats.

%While AI can help detect social engineering attacks, it does not mitigate the need for user training and awareness programs, quite the contrary, with AI-powered attacks the need for awareness and vigilance will likely grown even higher.

%Technologies such as deepfake (deep learning, fake) videos and live voice morphing will change the landscape of SE attacks.

% Tiivistelmäteksti on lyhyt, yleensä yhden kappaleen mittainen (maksimissaan noin 100 sanaa) selvitys kirjoituksen tärkeimmästä sisällöstä: mitä on tutkittu, miten on tutkittu ja mitä tuloksia on saatu.




\end{abstract}
\end{otherlanguage}


% Place ToC
%\newpage
\mytableofcontents

\mainmatter

%%%%%%%%%%%%%%%%%%%%%%%%%%%%%%%%%%%%%%%%%%%%%%%%%%%%%%%%%
%% STEP 4: Write the thesis.
%%%%%%%%%%%%%%%%%%%%%%%%%%%%%%%%%%%%%%%%%%%%%%%%%%%%%%%%%
%% Your actual text starts here. You shouldn't mess with the code above the line except
%% to change the parameters. Removing the abstract and ToC commands will mess up stuff.
%%
%% Command \include{file} includes the file of name file.tex.
%% A new page will be created at every \include command, 
%% which makes it appropriate to use it for large entities such as book chapters. Cannot be nested.
%% It is useful for a big project, as changing one of the include targets 
%% won't force the regeneration of the outputs of all the rest.
%% Alternatively, \input is a more lower level macro 
%% which simply inputs the content of the given file like it was copy&pasted there manually.



    %%%%%%%%%%%%%%%%%%%%%%%%%%%%%%%%%%%%%%%%%%%%%%%
    %% INTRODUCTION                              %%
    %%%%%%%%%%%%%%%%%%%%%%%%%%%%%%%%%%%%%%%%%%%%%%%


\chapter{Introduction\label{intro}}
\begin{comment}

Guides:
    - Find more things into this guides list
    - Page limit 1-2 pages
    - Context, who needs, what is needed, why is it a problem
    - Research question and research methodologies
    - Summary of results
    - What roles each section will have

TODO:
    [ ] Who is this thesis for
    [ ] Why should you read my thesis
    [ ] What is the research question and how it is answered
    [ ] How is this thesis organized, what is covered and what is deliberaly not covered and in what chapters (outline)

What to cover:
    - Motivate the user to read my thesis!
    - What is cybersecurity and why it's of paramount importance
    - What is SE?
    - Modern AI
        - Large Language Models
        - Emergence of ChatGPT and the like
    - Attacks
        - Deepfake synthetic content, videos, live voice morphing
        - Highly personalized phishing content (natural language processing)
        - Automated OSINT gathering
    - AI augments both attacks and countermeasures
    - Countermeasures
        - User awareness & training programs
        - Company policy & company culture
        - Real-time threat detection
        - Vulnerability detection
    - Emerging challenges
    - Motives for cybercrimes
        - Hard(er) to detect?
        - "Easy" wins?
    - Evolution of SE tactics and their impact on organizations
    - Intro to AI and its applications in cybersecurity
    - Research ga in the literature regarding the intersection of SE and AI


Info from lecture materials:
    - Johdannon tarkoituksena on kertoa yleiskielisesti työn tavoite. Kerrotaan (kuten tiivistelmässäkin, mutta laveammin), mitä on tutkittu, miten on tutkittu ja mitä tuloksia on saatu. Jotta kysymyksenasettelu ja tulokset on lukijan helppo oikein tulkita on syytä aloittaa johdanto asettelemalla tutkimus asiayhteyteensä, esimerkiksi kertomalla aluksi, minkälaisessa yhteydessä tarkasteluun otettavat haasteet esiintyvät ja keiden on ratkaisuista tarkoitus hyötyä.

    - Johdannon pituus määräytyy suhteessa koko kirjoitelman pituuteen. Parisivuinen kirjoitus ei erikseen otsikoitua johdantoa kaipaa, sillä se itsessään on laajennettu tiivistelmä. Kymmensivuisen kirjoituksen johdanto voi olla vaikkapa sivun tai puolentoista mittainen. Pro gradu -tutkielman 50-70-sivuiseen kokonaisuuteen tuntuu 2-4-sivuinen johdanto kohtuulliselta.

    - Johdanto kertoo siis lyhyessä, yleistajuisessa muodossa koko kirjoitelman kysymyksenasettelun, juonen sekä tulokset ja johtopäätelmät. Tämän luettuaan lukija voi päätellä, haluaako syventyä asiaan tarkemmin lukemalla koko kirjoituksen.

\end{comment}

The widespread adoption of information technology (IT) technologies and services has transformed every aspect of human life, from personal communication to business operations, and this reliance on devices and technologies is ever expanding. Although this digital revolution has opened up many opportunities, it has also brought about considerable vulnerabilities. One of the most dangerous threats to security and privacy is social engineering.

Social engineering (SE) is the art and science of manipulating people into performing actions or divulging confidential information. Rather than looking for technical vulnerabilities, SE relies on human interaction and exploits weaknesses in human psychology \citep{wangDefiningSocialEngineering2020}.

With the advent of modern artificial intelligence (AI), the landscape of social engineering is undergoing significant transformation, augmenting the sophistication and effectiveness of such attacks. Based on an in-depth literature review, this thesis explores the intersection of AI and SE, focusing on how contemporary AI technologies enhance the execution and impact of SE attacks and discusses the necessary actions to counter such advanced attacks. For illustration, case studies are also examined. Losses incurred by organizations in sophisticated social engineering attacks have topped \$100 million\footnote{https://www.justice.gov/usao-sdny/pr/lithuanian-man-pleads-guilty-wire-fraud-theft-over-100-million-fraudulent-business}.

Historically, SE relied heavily on human intuition and manual effort to deceive targets. Today, AI systems are more and more capable of automating and amplifying these deceptive practices, often with alarming precision. From sophisticated phishing schemes utilizing natural language processing to deepfake technologies creating convincing fake identities, AI-powered attacks represent a profound shift in the cybersecurity threat landscape.

Certain social engineering attacks that are of particular interest when it comes to advanced AI were chosen for more in-depth analysis, such as phishing, vishing (voice phishing), automated intelligence gathering, and deepfake-generated content, while leaving other attacks with less focus, such as dumpster diving, shoulder surfing, and tailgating.

Countermeasures that are especially covered are AI-based detection systems and how current cybersecurity training programs need to be augmented for the modern threat of AI.

%Next, key ideas and concepts in the field of SE are examined before moving on to the anatomy of a typical SE attack so that the actual attacks can be discussed in their proper context, all of this before we can finally delve into the countermeasures against the discsused attacks.

This thesis is organized into five chapters, each focused on a specific aspect of the interplay between AI and SE. The first chapters build the groundwork for further analysis, beginning with Chapter \ref{chapter:definition} which provides an overview of social engineering, including its definition and some background knowledge.

Chapter \ref{chapter:anatomy} delves into the various stages of a typical social engineering attack, highlighting the tactics and strategies employed by attackers. Chapter \ref{chapter:attacks} examines common attack methods and vectors, such as phishing, vishing (voice phishing), and pretexting, and explores how modern AI capabilities amplify their effectiveness.

Finally, with the groundwork being built, Chapter \ref{chapter:countermeasures} discusses AI-powered countermeasures and their potential to detect and prevent social engineering attacks along with other necessary steps to counter AI attacks. Finally, Chapter \ref{chapter:conclusions} concludes the thesis, summarizing the key findings and implications for the future of social engineering defense.



%Despite advancements in cybersecurity measures, SE remains a persistent threat. High-profile data breaches and corporate espionage incidents underscore the efficacy of social engineering techniques. Often, sophisticated firewalls, encryption algorithms and intrusion prevention systems (IPS) and intrusion detection systems (IDS) can be rendered useless by a cleverly worded email or an unsuspecting phone call. While it's often said that the most secure computer is one that is not plugged in, with SE an attacker could persuade someone to plug it in and turn it on.

%This study seeks to investigate the diverse aspects of SE, focusing on its techniques, effects, and prevention measures. Grasping the concept of SE is essential, not just for cybersecurity experts but also for people and organizations in every sector and at all employment levels. The increasing sophistication of social engineering attakcks coupled with the integration of artificial intelligence (AI) and technologies such as deepfakes (deep learning, fake) pose a growing risk that transcends traditional security frameworks. Because of this, extra attention is paid to how AI is and might contribute to the field of SE, both in terms of attacks and defenses.

%By dissecting real-world case studies and theoretical underpinnings, this thesis will provide a comprehensive overview of SE, contributing to the broader effort of mitigating its risks.

%Efficacy of SE training and awareness programs varies, with main challenges lying in user's interest and motivation towards the training and retention of learned data. Continued education, especially because SE is a constantly evolving field where the attackers never cease to improvise, is absolutely essential. Gamification, and users's inward motivation, go a long way in achievehing this.



%Usually when IT experts talk about an information system, they focus on the technical parts, the devices. But most, if not all, IT systems are useless without at least some interface with the people using it. Thus, especially from the SE standpoint, it's more beneficial to think of an information system as comprising the technology and its users. With this in mind, it becomes evident that the attack surface a hacker has is vast, and like water flowing down a river, the hacker often chooses the path of least resistance. Time and time again it has been proven that the easiest, most cost-effective way in is through the human factor of the IT system. Bolstering just one facet of this system is like fixing a leak or a potential leak only in some places, and not looking at the whole situation holistically. Security and privacy are a 360 degree, 24/7 thing.

%Currently, social engineering attacks are the biggest threat facing cybersecurity.


%This thesis is organized as follows. First, social engineering is defined and some history explained. Next, the different stages of a social engineering attack are examined. Then multiple common attack methods are discussed along with how modern AI powers them up. After that, AI-powered countermeasures are examined. The final chapter conclues this thesis.






    %%%%%%%%%%%%%%%%%%%%%%%%%%%%%%%%%%%%%%%%%%%%%%%
    %% DEFINITION                                %%
    %%%%%%%%%%%%%%%%%%%%%%%%%%%%%%%%%%%%%%%%%%%%%%%
    

\chapter{History and Definition of Social Engineering\label{definition}}

The term \textit{social engineering} is perhaps overused and is certainly misused, and for clarity, in this chapter a clearer definition is formed and some history reviewed. For the purpose of this thesis, we'll use the definition for SE given by \cite{wang_defining_2020}: "social engineeing is a type of attack wherein the attacker(s) exploit human vulnerabilities by means of social interaction to breach cybersecurity, with or without the use of techincal means and technical vulnerabilities."

The term \textit{social engineering} seems to have first appeared in an article titled "More on Trashing", which was published in September, 1984 on one of the earliest hacker magazines, The Hacker's Quarterly \citep{wang_defining_2020}, but the broader concepts of human exploitation date back much farther than this \citep{qin_investigation_2007}. For as long as humans have partaken in communication and trade, there have been those that have tried to exploit the system in their own favor, in an unethical and selfish way.

% impersonation
% phishing
% baiting
% tailgating

For as long as there has been social contracts, both written and oral or implied, there have been those that try to get away with cheating. Social engineering, thus, is not a new concept. In 19xx people who partook in actions akin to SE were called confidence men or con artists. They manipulated their victim to be complacent in the act.

As described by \cite{abiteboul}, the term \textit{social engineering} is perhaps overused and is certainly misused. What exactly constitutes social engineering? In this thesis, SE is defined as "social engineering is the deliberate act of convincing a victim, usually though the use of technology, to perform an action that may or may not be in their best interest". A definition of social engineering by Hadnagy has become quite popular "may or may not be in their best interest". HADNAGY

Some scholars have included acts like shoulder shurfing and dumpster diving as social engineering attacks \citep{abiteboul}, even though they do not rely on the manipulation of individuals.  However, since they don't fall neatly into the typical category of "hacking attacks", they fall into a gray area in between. Since user education and awareness programs should include training the user against discarding important documents without shredding them first or by putting them into boxes designated "secure documents", and people need to be on the lookout for people gazing over their shoulders, especially when entering sensitive information such as usernames, passwords and other access codes, they are included in this thesis.

It's also necessary to go over some basic terminology related to the field of SE.

Two common attack methods, namely \textbf{dumpster diving} and \textbf{shoulder surfing}, are often categorized as social engineering attacks, but do not necessarily fall under the SE category, as they do not include direct social interaction with the victim \citep{wang_defining_2020}. In other words, the compliance of the victim is not necessary in these two types of attacks. However, since they are non-technical in that no devices or technologies need to be used, they fall in a "gray area". They require neither the use of devices or technologies nor the psychological manipulation of the user. Thus, it has been easy to refer to them as SE attacks. In this paper, they are still discussed because SE training and awareness programs should include mentioning about these.








        %%%%%%%%%%%%%%%%%%%%%%%%%%%%%%%%%%%%%%%%%%%%%%%
        %% Open-Source Intelligence (OSINT)          %%
        %%%%%%%%%%%%%%%%%%%%%%%%%%%%%%%%%%%%%%%%%%%%%%%

\section{Open-Source Intelligence, OSINT}

OSINT, sometimes written as OS-INT, means open-source intelligence. Like the name implies, it involves gathering of intelligence data from publicly locatable sources, such as from the target company's website, or from the social networking profiles of an individual or from other public records.

As people have adopted to using social media as sometimes their primary means of communication, a lot of exploitable data is shared on these platforms as well.
\chapter{Stages of Social Engineering Attacks\label{stages}}

Social engineering attacks are deployed in multiple stages. Various scholars define these stages a little differently, but they generally follow a pattern such as the one defined by \cite{wang_defining_2020}:

\begin{enumerate}
    \item Collect information about the target
    \item Develop a relationship with the target
    \item Exploit the available information and the built relationship by executing the attack
    \item Exit without leaving or having cleaned any traces
\end{enumerate}

SE attacks can be detected, but they cannot be fully stopped \citep{wang_defining_2020}.

Next, each is these four stages is examined in detail.

\section{Intelligence Gathering \& Open-Source Intelligence}

First stage is research or information collection. This relies on what's known as OSINT, or open-source intelligence. OSINT is any information that is publicly available via the Internet or other means and that doesn't require any breach of security to be accessed.

\section{Developing the Relationship}

\section{Exploiting the Victim}

\section{The Exit}

Hadnagy et al. emphasize that in whatever you do, leave your target better off for having met with. This is in stark contrast to typical SE attacks.

Any bridges should not be burned but the relationship should be exploitable in the future as well


    %%%%%%%%%%%%%%%%%%%%%%%%%%%%%%%%%%%%%%%%%%%%%%%
    %% ATTACK METHODS                            %%
    %%%%%%%%%%%%%%%%%%%%%%%%%%%%%%%%%%%%%%%%%%%%%%%

\chapter{AI-Powered Attacks\label{chapter:attacks}}
\begin{comment}

Guides:
    - About 3-4 pages

TODO:
    [ ] 

What to cover:
    - Attacks
        - Deepfake generated synthetic media
            - Videos
            - Images
            - Audio
            - Real-time voice morphing
    
Literature:
    - 

\end{comment}

Awareness of various social engineering (SE) attacks is crucial for professionals across all industries, not just those in cybersecurity. This chapter provides an overview of the most common SE attack methods. 

%A later chapter will examine how modern AI augments these attacks, and how users of information systems need to be trained to counter these new, advanced threats.

Some attacks that are discussed here are not always considered a type of SE attack, specifically shoulder surfing and dumpster diving, as these do not require the co-operation of the victim \citep{wangDefiningSocialEngineering2020}. However, since they are often used in conjunction with other SE attacks, and since training for against them is often included in SE training and awareness programs, they are explained here

%and because attackers are constantly improvising their attack methods, some manipulation used in conjunction with these methods, such as convincing employees that the company will destroy any documents put in the general waste in a safe way, may be used.

To better understand the threat posed by SE, it is essential to examine the diverse strategies employed by attackers. The following are some of the most common and effective SE attack methods. We'll also analyze how the emergence of modern AI technologies might, or already has, powered up these types of attacks.

The field of social engineering contains a plethora of attacks, and for the purposes of this thesis this selection had to be narrowed down. Due to their relevance, phishing it its various forms, deepfake-generated content and automated OSINT were chosen for a more in-depth analysis.

Countermeasures against these attacks are examined in Chapter \ref{chapter:countermeasures}.









        %%%%%%%%%%%%%%%%%%%%%%%%%%%%%%%%%%%%%%%%%%%%%%%
        %% Phishing                                  %%
        %%%%%%%%%%%%%%%%%%%%%%%%%%%%%%%%%%%%%%%%%%%%%%%

\section{Phishing}
\begin{comment}
    
    - 

\end{comment}


As the quintessential SE attack, \textbf{phishing} is characterized by malicious attempts to gain sensitive information from unaware users, usually via email and by using spoofed websites that look like their authentic counterparts. Phishing has been around since 1996, when cybercriminals began using deceptive emails and websites to steal AOL (America Online) account information from gullible users \citep{wangDefiningSocialEngineering2020}.

Table \ref{tab:placeholder_label} presents a sample, legitimate URL and a non-functional URL used for phishing purposes. The target user is hoped by the attacker to be unaware of the different domains these seemingly similar URL's redirect the user. The phishing URL is taken from a dataset of known phishing websites but its domain is changed to IANA's example.com domain for security. The legitimate URL is from PayPal's website and is functional as of writing.

\begin{table}[h]
    \centering
    \begin{tabular}{|c|c|}
        \hline
        Legitimate & Phishing \\ \hline
        paypal.com/us/signin &paypal.com.cgi-bin.788a5.example.com/ \\ \hline
    \end{tabular}
    \caption{Example of a legitimate URL and a non-working phishing URL.}
    \label{tab:placeholder_label}
\end{table}

\textbf{Spear phishing} is a more targeted version of phishing, where attackers customize their deceptive emails to a target individual or organization. Unlike with generic phishing attempts, this type of phishing involves gathering detailed information about the victim, via OSINT or otherwise, such as their name, position and contacts to craft a convincing and personalized message \citep{salahdine_social_2019}. This tailored approach increases the likelihood of the victim falling for the scam.

Last on the list of phishing attacks is \textbf{whaling}. Whaling, also known as CEO fraud, is a highly targeted phishing attack aimed at high-profile individuals within an organization, such as executives or senior management, "the big whales" \citep{abraham_overview_2010}. The attackers carefully research their targets to create convincing and typically urgent messages that appear to come from trusted sources, often impersonating colleagues, business partners, or government agencies. The goal is often to authorize large financial transactions or to leverage the target's authority and access within the company.

Two additional types of phishing need to be addressed, and they are \textbf{vishing} or voice phishing and \textbf{smishing} or SMS (simple message system) phishing. Despite having complicated names, the idea behind them are quite simple.

Regular phishing doesn't usually require OSINT but spear phishing does.

        %%%%%%%%%%%%%%%%%%%%%%%%%%%%%%%%%%%%%%%%%%%%%%%
        %% Deepfake-augmented                        %%
        %%%%%%%%%%%%%%%%%%%%%%%%%%%%%%%%%%%%%%%%%%%%%%%

%\section{Deepfake-augmented attacks}
\begin{comment}
    
    -

\end{comment}



%\section{Automated OSINT}
\begin{comment}
    
    - 

\end{comment}


        %%%%%%%%%%%%%%%%%%%%%%%%%%%%%%%%%%%%%%%%%%%%%%%
        %% Other attack methods                      %%
        %%%%%%%%%%%%%%%%%%%%%%%%%%%%%%%%%%%%%%%%%%%%%%%

\section{Other attack methods}
\begin{comment}
    
    - 

\end{comment}


This section covers other social engineering attacks that were not chosen for more in-depth analysis due to their weaker affinity to be amplified by modern AI. These attacks include tailgating, shoulder surfing, and dumpster diving. However, advances in related fields, such as robotics \citep{postnikoffRobotSocialEngineering2018}, could make these attacks more relevant in the future. All of them can be used as part of a social engineering attack chain, even if their use would not rely on AI directly, and this is why it is important that they be mentioned in this thesis.





%%\section{Baiting}

%%\textbf{Baiting} is a similar attack method to phishing, discussed above, but emphasizes luring the victim via enticement strategies \citep{conteh_cybersecurityrisks_2016, salahdine_social_2019}. This technique exploits the target's curiosity or greed to gain unauthorized access to resources or premises or to obtain sensitive information.

%%A common type of a baiting attack is what's known as a "USB drop" where the attacker loads malware or Trojan software into USB drives that are often tagged with an enticing title such as "Layoff Plans 2024" or "Sarah's private photos" and which are subsequently dropped at convenient locations for employees of the target company to find. These locations could be easily accessed, such as the company's vicinities, or more hard-to-reach places such as a protected parking lot or even employee restrooms or dining areas. A famous experiment in which 297 flash drives were dropped on university campuses concluded that the success rate (drives connected) was between 45 \% -- 98 \%, with the first drive connected within less than six minutes \cite{tischer_users_2016}.

%%A USB baiting experiment showed that 15 out of 20 USB drives were found by employees and all were plugged in to the company's computers \citep{wang_social_2021}.

%%OSINT can be useful here as well, with a bit of searching the tag that is attached to the thumbdrive could have a more interesting and suitable title. Attackers can also easily order USB thumbdrives, as an example, with custom logos printed on them from various online retailers for increased effect. The logo could even be from a competing company's one, if there's been heated rivalry between the two.

%% https://www.darkreading.com/perimeter/social-engineering-the-usb-way

%% online offers, free downloads, gifts, fake job posting
%% plays on human trust, safe looking items, greed, curiosity
%% fake software updates, malicious ads

%% seuraukset: malware, trojan, data breach, financial loss, identity theft
%% red flags? unsolicitied offers, too good to be true deals, unfamiliar sources











        

%%Situations can change quickly, and the attacker must be quick on their feet to respond to these changes. Training in multiple possible ways the target might respond or react is highly beneficial, such as threatening to call the police or the IT support.

%%A high level of confidence in the pretext and the pretexted role is often cited to be a necessity, however, someone may very well engage a pretext as a novice, nervous new employee, which suggests that rather than stating a high level of (outwardly visible) confidence as a neceissty, the better wording would be a suitable level of confidence. In all cases, the attacker must believe their own pretext and act out the attack like an actor in a movie, to "become the pretext".

%% obtain sensitive info, access to systems or locations
%% must be believable and relevant, should fit with the target's environment and expected interactions
%% kerro kuinka exploitatataan luottamusta
%% kommunikaatiotaidot, professionalism, itsevarmuus (sopiva taso, jos esim. pretextaa aloittelijaa itsevarmuus voikin olla alhaisempi)
%% tyypillisiä skenaarioita: co-worker, kuljetushtiö, security audit, IT tuki, law enforcement, 
%% mitä dataa haetaan: henk koht tiedot, pankkitiedot, kirjautumistiedot, joka päivä vaihtuva "PIN" koodi, insider info
%% suojautuminen: soittajan id vahvistus (varsinkin tuntemattomien soittajien), caller spoofing vahvistus, mitä tietoja saa antaa, kenelle ja millä edellytyksillä, 
%% pretexting on laitonta monessa paikkakunnassa, paikassa, SE white hat audit
%% oikean maailman incidentit












        %%%%%%%%%%%%%%%%%%%%%%%%%%%%%%%%%%%%%%%%%%%%%%%
        %% Tailgating                                %%
        %%%%%%%%%%%%%%%%%%%%%%%%%%%%%%%%%%%%%%%%%%%%%%%

\begin{comment}
    
    - 

\end{comment}

\textbf{Tailgating}, also known as \textbf{piggybacking}, is a social engineering tactic that involves following an authorized person through an access-controlled passage, such as a security gate. This type of attack exploits individuals with temporary access rights, such as delivery personnel or maintenance workers \citep{conteh_cybersecurityrisks_2016}. The attacker may use manipulation to gain access, for instance, by carrying a heavy object and asking for assistance or pretending to be a delivery person with a forged pretext.

%%The attacker may have found out through OSINT that the company is expecting a delivery and then don the suit of a delivery company and escort people going in from their smoking/coffee breaks outside. Another type is where the attacker has forged a pretext where he explains that he had "forgotten" his security ID inside and his boss is already very angry with him and that he's afraid he's going to lose his job if he keeps messing up like this, playing on people's sympathy.

%%In another scenario, the attasacker may use psychological manipulation to evoke sympathy, claiming to have forgotten their security ID and worrying about losing their job. This tactic preys on people's natural instinct to help and be polite. The mundane routine of passing through security gates can also lead to a false sense of security, making individuals less vigilant.

%%This type of attack exploits people's natural tendency to want to help and be polite. Going in through security gates is a daily chore for many, which due to having become a mundane activity, has lessened the individual's vigilance.

%% may use machine learning algorithms to identify vulnerabilites in access control systems, making tailgating attacks even more dangerous












        %%%%%%%%%%%%%%%%%%%%%%%%%%%%%%%%%%%%%%%%%%%%%%%
        %% Shoulder Surfing                          %%
        %%%%%%%%%%%%%%%%%%%%%%%%%%%%%%%%%%%%%%%%%%%%%%%

%%\section{Shoulder Surfing}

%% STATUS: almost complete, need sources

Observing people enter sensitive information, such as login details or financial data, without their knowledge or approval is called \textbf{shoulder surfing}. The name implies a person watching over the shoulder of someone when they are typing or viewing sensitive data, but the attack surface is actually far larger. The attacker could use cameras, either hacked, or those placed there by the attacker for this purpose, or even everyday objects such as binoculars.

%%The proliferation of high-resolution cameras, such as those with HD or 4K resolution, has exacerbated the issue of unauthorized information gathering since security cameras of the past didn't have the resolution necessary to show what's being viewed on a screen. Going through tens or even hundreds of hours of video material where sensitive information might be visible on an individual's or an employee's screen is time-consuming and tedious, but with the help of AI this job can be carried out with more ease. AI can not only turn words in a video into text, but it can also summarize the found text and search for anything that could be exploited.

%% vigilance in public places
%% HD and 4K cameras
%% social norms, vigilance
%% pankkiautomaatti
%% kohteet PIN, salasanat, turvakoodit, henk koht tiedot
%% puhelimet, draw a symbol to access, halvat vakoilulaitteet aliexpress
%% privacy screens, HP SureView
%% toimiston layout joka hankaloittaa shoulder surfing
%% etätyö













        %%%%%%%%%%%%%%%%%%%%%%%%%%%%%%%%%%%%%%%%%%%%%%%
        %% Dumpster Diving                           %%
        %%%%%%%%%%%%%%%%%%%%%%%%%%%%%%%%%%%%%%%%%%%%%%%

%%\section{Dumpster Diving}

%% STATUS: lähes valmis

As the name implies, \textbf{dumpster diving} refers to the practice of going through an organization's or an individual's trash in order find sensitive information that should have been disposed of properly \citep{syafitri_social_2022}. Rummaged content may yield interesting results, typically in the form of documents such as papers but also media devices such as optical discs, hard drives or USB thumbdrives, to be used as-is or as part of a future SE attack. Since the focus of this thesis is on AI and not robotics, dumpster diving is not examined further.

%%Since dumpster diving does not include manipulation of people in its purest form, rather relying on the improper care of documents and other material, it is sometimes not classified as a social engineering attack \citep{wang_defining_2020}. However, it is often considered a precursor or supplementary tactic in broader social engineering schemes, as the information gathered can be used to craft more convincing and targeted attacks.

%% However, one could imagine scenarios where manipulation of people is used as part of a dumpster diving attack, for example convincing people that documents put into regular waste bins will be shredded as part of a new company policy, or that all discs and USB drives are constantly being encrypted and thus they can be safely discarded in the regular trash.









%% Käsittelylukujen työnjako määräytyy käsiteltävän asian luonteen mukaisesti. Lukijan oh￾jailemiseksi kukin pääluku kannattaa aloittaa lyhyellä kappaleella, joka paljastaa mikä kyseisen luvun keskeisin sisältö on ja kuinka aliluvuissa asiaa kehitellään eteenpäin. Er￾ityisesti kannattaa kiinnittää huomiota siihen, että lukijalle ilmaistaan selkeästi miksi kutakin asiaa käsitellään ja miten käsiteltävät asiat suhtautuvat toisiinsa. Jäsentelyongelmista kielivät tilanteet, joissa alilukuja on vain yksi, tai joissa käytetään useampaa kuin kahta tasoa (pääluku ja sen aliluvut). Kolmitasoisia otsikointeja saate￾taan tarvita joissakin teknisissä dokumenteissa perustellusti, mutta nämä muodostavat poikkeuksen.

%% Perusohjeena on käyttää tekstin rakenteellisesti painokkaita paikkoja, kuten lukujen avauk￾sia ja teksikappaleiden aloitusvirkkeitä juonenkuljetukseen ja informaatioaskeleiden sito￾miseen toisiinsa. Tekstikappaleiden keskiosat, samoin kuin lukujen keskiosat selostavat asiaa vähemmän tuntevalle yksityiskohtia, kun taas aihepiirissä jo sisällä olevat lukijat voivat alkuvirkkeitä silmäilemällä edetä tekstissä tehokkaasti eksymättä tarinan juonesta.

%% Kullakin kirjoittajalla on oma temponsa, joka välittyy lukijalle tekstikappaleiden pitu￾udessa ja niihin sisällytettyjen ajatuskulkujen mutkikkuudessa. Kussakin tekstikappaleessa pitäisi pitäytyä vain yhdessä informaatioaskelessa tai olennaisessa päättelyaskelessa, muuten juonen seuraaminen käy raskaaksi olennaisten lauseiden etsiskelyksi. Yksivirkkeisiä tek￾stikappaleita on syytä varoa.

\textbf{Reverse social engineering} is a special attack vector wherein the attacker advertises himself as providing a solution to problems within a certain domain, usually in the IT/tech sphere, then artificially creates a problem and waits for the victim to contact them.

Table \ref{tab:attacks} lists some of the most common SE attacks, both in their old-school, pre-AI era and as AI-powered to give a clear overview.

\begin{table}[ht]
    \centering
    \begin{tabularx}{\textwidth}{X X X}
    %    \begin{tabular}{m{3cm}m{5,5cm}m{5,5cm}}

        \toprule
        \textbf{Attack Type} & \textbf{Pre-AI} & \textbf{AI-Powered} \\
        \midrule
        Phishing \& Smishing & Generic messages, low personalization, easily detectable &  Highly personalized, uses NLP, adaptive, real-time learning from responses \\
        \midrule
        Spear Phishing & Manual research for personalization, time consuming & Automated, deep-learning based personalization, rapid deployment \\
        \midrule
        Vishing & Human-driven calls, script based & AI-generated voice, dynamic script adjustments based on real-time conversation analysis \\
%        \midrule
        %Baiting & Physical media (e.g. USB drives), opportunistic & AI-driven malware distribution through personalized and enticing digital baits \\
        \midrule
        Impersonation & Human-driven calls, script based & AI-generated voice, dynamic script adjustments based on real-time conversation analysis \\
        \bottomrule
        
    \end{tabularx}
    %\end{tabular}
    \caption{Comparison of Social Engineering Attack Methods (Pre and Post-AI)}
    \label{tab:attacks}
\end{table}


    %%%%%%%%%%%%%%%%%%%%%%%%%%%%%%%%%%%%%%%%%%%%%%%
    %% COUNTERMEASURES                           %%
    %%%%%%%%%%%%%%%%%%%%%%%%%%%%%%%%%%%%%%%%%%%%%%%


\chapter{Countermeasures\label{chapter:countermeasures}}
\begin{comment}


\end{comment}

In this chapter, AI-powered countermeasures against both classic social engineering and AI-powered social engineering are examined. This chapter is divided into two parts: tech-oriented countermeasures such as deepfake detection mechanisms and human-oriented countermeasures such as training and awareness programs. Tech-oriented countermeasures are examined first since the human-oriented measures rely and build upon them. The division is not always clear cut and is made only to simplify the reading experience.

Finally, Chapter \ref{chapter:conclusions} concludes the thesis.

        %%%%%%%%%%%%%%%%%%%%%%%%%%%%%%%%%%%%%%%%%%%%%%%
        %% Technology oriented                       %%
        %%%%%%%%%%%%%%%%%%%%%%%%%%%%%%%%%%%%%%%%%%%%%%%

\section{Technology oriented}
\begin{comment}
    
    - 

\end{comment}

Countermeasures against AI-powered threats to SE that fall into the technology-oriented, versus human-oriented, are discussed.

They include automated phishing email detection powered by AI and machine learning algorithms.


        %%%%%%%%%%%%%%%%%%%%%%%%%%%%%%%%%%%%%%%%%%%%%%%
        %% Human oriented                            %%
        %%%%%%%%%%%%%%%%%%%%%%%%%%%%%%%%%%%%%%%%%%%%%%%


\section{Human oriented}
\begin{comment}
    
    - The best defense against SE attacks is an educated, conscious user
    - User education should be continuous and not a one-off event

\end{comment}

Regular and comprehensive training programs are vital to educate employees about social engineering tactics.

Conducting concensual, simulated phishing attack campaigns, via numerous channels such as email, SMS and even phone/VoIP, allows organizations to assess the suspectibility of their employees to SE tactics. These exercises help identify vulnerabilities in the workforce, enabling targeted training and reinforcing the importance of scrutiziing unsolicited communication. Feedback from these simulations can be a powerful tool for personnel development. It's important that Employees should be informed beforehand that such campaigns may be intermittently run, which has the double benefit of keeping the employees on their guard and also not causing unneccessary bad emotions from "being tricked" by their own company \citep{hadnagySocialEngineering2018, mitnickArtDeceptionControlling2003}.

Regarding bating attacks, company policy should enforce antivirus scanning of attached media, such as USB thumbdrives, before they can be accessed. Personnel should be instructed to never pick up and plug in found media \citep{salahdine_social_2019}. However, if media is found that is suspected of being used as part of a baiting attack, the employee could notify the front desk or the IT services that such a media device was found within the premises, rather than just ignoring it 

Because AI can source the Internet automatically for OSINT, it's imperative for people to know to be careful of what they share, with whom and when. There's a difference between sharing that one is going to a company picnic next week, and posting videos about a picnic that had taken place 3 days prior.



In this chapter, some key countermeasures against pre-AI social engineering as well as AI-powered SE were looked at. The final chapter \ref{conclusions} ties everything together and provides speculation about the current state of the AI SE landscape as well as future prospects.




    %%%%%%%%%%%%%%%%%%%%%%%%%%%%%%%%%%%%%%%%%%%%%%%
    %% CONCLUSIONS                               %%
    %%%%%%%%%%%%%%%%%%%%%%%%%%%%%%%%%%%%%%%%%%%%%%%


\chapter{Conclusions\label{conclusions}}
\begin{comment}
    
Guides:
    - As many pages as it takes (thesis max is around 20 pages)

TODO:
    [ ] How AI has augmented SE attacks and countermeasures
    [ ] Gap in the literature

What to cover:
    - How AI has augmented SE attacks and countermeasures
    - Gap in the literature regarding SE and AI intersection?
    - Analysis on where AI-powered SE attacks might be headed in the future
        - Also about robotics and human-like actors
    - What organizations and individuals need to do regarding the evolving landscape of SE attacks

Speculation:
    - Drones dropping USB thumbdrives?
    - Human-like android as threat actors
    - Impact of robotics on dumpster diving, shoulder surfing and baiting
    
Literature:
    - Gen and detection of deepfakes

From training material:
    - Yhteenveto vaatimattomimmillaan on vain lyhyt kertaus kirjoituksen keskeisistä asioista. Arvokkaamman yhteenvedon saa aikaan kommentoimalla työn tulosten arvoa, työn liittymistä ympäristöön ja tulevaisuudennäkymiä. Tällaiset arviot huolellisesti perusteltava.

\end{comment}


As we've seen, SE is still as much a threat as it has ever been, despite major efforts to the contrary.

What's certain is that we can count on AI developing, AI-based social engineering attacks evolving with it, and the need for continuous, innovative user training growing in the future. Attackers and defenders are playing a never-ending game of "cat \& mouse" where nobody can rest.

X in Y references that training users effects will wear off in 3 weeks, necessiting continous retraining approaches.

I'll end with the question that I started with; what, if anything, can the end-user trust anymore? And perhaps, with the advances in AI technology, the answer is "no-one".


%Modern AI can assist with the deployment of the pretext by being a sparring partner to the attacker, giving a safe "sandbox" to try out potential attacks and find ways in which the target might react.

%%%%%%%%%%%%%%%%%%%%%%%%%%%%%%%%%%%%%%%%%%%%%%%%%%%%%%%%%
%\cleardoublepage                          %fixes the position of bibliography in bookmarks
%\phantomsection
\addcontentsline{toc}{chapter}{\bibname}  % This lines adds the bibliography to the ToC
\printbibliography

%%%%%%%%%%%%%%%%%%%%%%%%%%%%%%%%%%%%%%%%%%%%%%%%%%%%%%%%%
\backmatter
\begin{appendices}

%% A sample Appendix
%%
\appendix{Sample Appendix\label{appendix:sample}}

You can add one or more appendices to your thesis.




%% another appendix
%%
\appendix{Instructions for LaTex}

\section{General Setup}

In the HY-CS-main.tex file you will find the following STEPS 0--5. Below you can find related instructions.
\vspace{0.5cm}

\textbf{STEP 0 -- Access the thesis template}

\begin{itemize}
\item Import the thesis template into a new Overleaf project. The easiest way to do it is to:
\begin{itemize}
    \item Obtain a zip file of the LaTeX template from the webpage of your programme.
    \item Go to \url{https://www.overleaf.com/edu/helsinki} and login to Overleaf with your university credentials.
    \item Go to the list of your projects at \url{https://www.overleaf.com/project}, click ``New Project'' and ``Upload Project''.,  the projects under your account 
    \item Then upload the zip with the template.
    \item You are now ready to write your thesis in Overleaf by editing the template, you can start by renaming the project.
\end{itemize}
\end{itemize}


{\textbf{STEP 1 -- BSc or MSc thesis?}}
\begin{enumerate}
\item Select whether your are writing BSc (tkt) or MSc (csm for CS) thesis.
\item Select your language: \texttt{finnish}, \texttt{english}, or \texttt{swedish}.
\item If you are writing MSc select your line / track.
\end{enumerate}


{\textbf{STEP 2 -- Set up your personal information}}

\begin{enumerate}
\item Specify the title of your thesis with \texttt{\textbackslash title\{\}}.
\item Specify your name to the author field with \texttt{\textbackslash author\{\}}.
\item Specify the names of your supervisors of the thesis with \texttt{\textbackslash supervisors\{\}}.
\item Specify the keywords of the thesis with \texttt{\textbackslash keywords\{\}}.
\item Specify the ACM classification terms of the thesis with \texttt{\textbackslash classification\{\}}. See \url{https://dl.acm.org/ccs} for more information.
\end{enumerate}

{\textbf{STEP 3 -- Write your abstract}}

\begin{itemize}
\item You can have the abstract in multiple languages with the \texttt{otherlanguages} environment. The example below shows how to provide an English abstract: 

\begin{verbatim}
\begin{otherlanguage}{english} 
\begin{abstract}
Your abstract text goes here. 
\end{abstract} 
\end{otherlanguage}
\end{verbatim}

\end{itemize}

{\textbf{STEP 4 -- Writing your thesis}}

\begin{enumerate}
\item There are some minimal contents and instructions below 
\item Remove, or comment out, this appendix from your thesis.
\end{enumerate}

{\textbf{STEP 5 -- Set your bibliography style}}

\begin{itemize}
\item The default is Author-Year style (Einstein, 1905), but it can be easily changed to numbered [1] or alphabetical [Ein05] , as the examples of these are in comments.
\item Discuss the style to use with your supervisor.
\end{itemize}

\section{Bibliography in Latex}

The bibliography is defined in a separate \texttt{.bib} file. For this template, it is named \texttt{bibliography.bib} and includes the content show in Figure~\ref{bibexamples}.

Chapter Bibliography lists all the works that you refer to in your text. You refer to the works in the bibliography using an appropriate \emph{citation key}.
%
%This thesis template contains an example of a bibliography.


References are done using \texttt{\textbackslash citep\{einstein\}}, which generates in text a citation formatted according to the selected style \citep{einstein}, or \texttt{\textbackslash citep\{latexcompanion,knuth99\}}, which generates \citep{latexcompanion,knuth99}. 
As examples of a different kinds of citations (see how these look in the Latex source), we can write \citep{einstein} to refer to the work written by \citeauthor{einstein} in \citeyear{einstein}, because the work by \citet{einstein} appears in the bilbliography included in this template.

Note that there are different possible styles for the bibliography and citation keys.
%
Consult your supervisors on the chosen style -- and once you arrive at a preferred style, use it consistently throughout the thesis.

\begin{figure}[ht]
    \centering
    \begin{scriptsize}
\begin{verbatim}
@article{einstein,
    author =       "Albert Einstein",
    title =        "{Zur Elektrodynamik bewegter K{\"o}rper}. ({German})
        [{On} the electrodynamics of moving bodies]",
    journal =      "Annalen der Physik",
    volume =       "322",
    number =       "10",
    pages =        "891--921",
    year =         "1905",
    DOI =          "http://dx.doi.org/10.1002/andp.19053221004"
}
 
@book{latexcompanion,
    author    = "Michel Goossens and Frank Mittelbach and Alexander Samarin",
    title     = "The \LaTeX\ Companion",
    year      = "1993",
    publisher = "Addison-Wesley",
    address   = "Reading, Massachusetts"
}

@book{knuth99,
    author    = "Donald E. Knuth",
    title     = "Digital Typography",
    year      = "1999",
    publisher = "The Center for the Study of Language and Information",
    series    = "CLSI Lecture Notes (78)"
}\end{verbatim}
\end{scriptsize}
    \caption{Examples of bibliographic reference in .bib file.}
    \label{bibexamples}
\end{figure}

%In the last reference url field the code \verb+%7E+ will translate into \verb+~+ once clicked in the final pdf.

\section{Some instructions about writing in Latex}

The following gives some superficial instructions for using this template for a Master's thesis. For guidelines on thesis writing you can consult various sources, such as university courses on scientific writing or your supervisors.

For more detailed instructions, just google, e.g., "Overleaf table positioning", and your chances of finding good info are pretty good.  


\section{Figures}
Besides text, here are simple examples how you can add figures and tables in your thesis.
Remember always to refer to each figure in the main text and provide them with a descriptive caption.

Figure~\ref{fig:logo} is an example of a figure in the document (see the source about how to add them). 
%Using figures is particularly useful to display plots of experimental results.

\begin{figure}[ht] 
\begin{center}
\includegraphics[width=0.3\textwidth]{template/figures/HY-logo-ml.png}
\caption{University of Helsinki flame-logo for Faculty of Science.\label{fig:logo}}
\end{center}
\end{figure}

\section{Tables}

Table~\ref{table:results} gives an example of a table.
Remember always to cite the table in the main text, table captions go on top of the table. 

\begin{table}[h] % h positions the table here, t! would force on top of the page, or example.
\begin{center}
\caption{Experimental results.\label{table:results}} % caption is here to make it on top
\begin{tabular}{l||l c r} 
Experiment & 1 & 2 & 3 \\ 
\hline \hline 
$A$ & 2.5 & 4.7 & -11 \\
$B$ & 8.0 & -3.7 & 12.6 \\
$A+B$ & 10.5 & 1.0 & 1.6 \\
\hline
%
\end{tabular}
\end{center}
\end{table}



%% yet another appendix
%%
\appendix{Tutkielmapohjan käyttöohjeet}
\label{appendix:instructions_finnish}

\section{Ensiaslkeleet}

\texttt{HY-CS-main.tex} tiedosto sisältää viisi askelta STEPS 0--5. Alla on kuvattu, mitä nämä askeleet tarkoittavat ja miten niitä seuraamalla luot pohjan tutkielmallesi.
\vspace{0.5cm}

\textbf{STEP 0 -- Kopioi tutkielmapohja}

\begin{itemize}
\item Hae tutkielmapohja uuteen Overleaf-projektiin. Tämä käy helpoiten seuraavasti:
\begin{itemize}
    \item Lataa Latex-pohjan zip-tiedosto koulutusohjelman sivuilta.
    \item Mene osoitteeseen \url{www.overleaf.com/edu/helsinki} ja kirjaudu Overleafiin yli\-opiston tunnuksillasi.
    \item Overleafissa (\url{https://www.overleaf.com/project}), klikkaa ``New Project'' and ``Upload Project''.
    \item Valitse lataamasi tutkielmapohjan zip-tiedosto.
    \item Nyt voit lähteä kirjoittamaan tutkielmaasi suoraan pohjaan, voit aloittaa esim. vaihtamalla projektin nimen.
\end{itemize}
\end{itemize}


{\textbf{STEP 1 -- BSc vai MSc tutkielma?}}
\begin{enumerate}
\item Valitse (tiedostossa \texttt{HY-CS-main.tex}) oletko tekemässä BSc (tkt) vai MSc (csm tietojenkäsittely) tutkielmaa.
\item Valitse kieli jolla kirjoitat tutkielman: \texttt{finnish}, \texttt{english} tai \texttt{swedish}.
\item Jos olet kirjoittamassa maisterintutkielmaa, valitse linja/opintosuunta.
\end{enumerate}


{\textbf{STEP 2 -- Aseta henkilökohtaiset tietosi}}

\begin{enumerate}
\item Kirjoita alustava otsikko tutkielmallesi: \texttt{\textbackslash title\{\}}.
\item Kirjoita oma nimesi kohtaan \texttt{\textbackslash author\{\}}.
\item Lisää ohjaajien nimet \texttt{\textbackslash supervisors\{\}}.
\item Määrittele avainsanat \texttt{\textbackslash keywords\{\}}.
\item Määritä tutkielmasi ACM luokittelutermit \texttt{\textbackslash classification\{\}}. Ks. lisätietoa: \url{https://dl.acm.org/ccs}.
\end{enumerate}

{\textbf{STEP 3 -- Kirjoita tiivistelmä}}

%\begin{itemize}
%\item 
Voit kirjoittaa tiivistelmän (koko tiivistelmäsivu) eri kielillä \texttt{otherlanguages}-ym\-pä\-ris\-tön avulla. Alla esimerkki jolla kirjoitat englanninkielisen tiivistelmän muulla kuin englannin kielellä kirjoitettuun tutkielmaan:

\begin{verbatim}
\begin{otherlanguage}{english} 
\begin{abstract}
Your abstract text goes here. 
\end{abstract} 
\end{otherlanguage}
\end{verbatim}

%\end{itemize}

{\textbf{STEP 4 -- Kirjoita tutkielma}}

\begin{enumerate}
\item Kirjoittamisesta Latexilla löydät hieman ohjeita alempaa.
\item Poista tämä liite ja muu ohjeistus tutkielmastasi, esim. kommentoimalla.
\end{enumerate}

{\textbf{STEP 5 -- Aseta kirjallisuuslähdeluettelon tyyli}}

\begin{itemize}
\item Oletustyylin tekijä-vuosi, eli (Einstein, 1905), voit vaihtaa viittaustyylin (tiedostossa \texttt{HY-CS-main.tex}) helposti (eri mallit kommentoituna) esim. numeroituun [1], tai aakkostyyliin [Ein05].
Lisää ohjeita liittyen viittaustyylin säätämiseen {Bib}\TeX issä löytyy verkosta: \url{https://ctan.org/pkg/biblatex}
\item Sovi käytettävä tyyli ohjaajasi kanssa. 
\end{itemize}

\section{Kirjallisuusviitteet Latexissa}

Kirjallisuuslähteet ylläpidetään erillisessä .bib-tiedostossa. Tässä tutkielmapohjassa käy\-te\-tyt kirjallisuuslähteet, joista esimerkkejä kuvassa~\ref{bibexamples-fi}, löytyvät tiedostosta\newline \texttt{bibliography.bib}.

\begin{figure}[ht]
    \centering
    \begin{scriptsize}
\begin{verbatim}

@article{einstein,
    author =       "Albert Einstein",
    title =        "{Zur Elektrodynamik bewegter K{\"o}rper}. ({German})
        [{On} the electrodynamics of moving bodies]",
    journal =      "Annalen der Physik",
    volume =       "322",
    number =       "10",
    pages =        "891--921",
    year =         "1905",
    DOI =          "http://dx.doi.org/10.1002/andp.19053221004"
}

@book{latexcompanion,
    author    = "Michel Goossens and Frank Mittelbach and Alexander Samarin",
    title     = "The \LaTeX\ Companion",
    year      = "1993",
    publisher = "Addison-Wesley",
    address   = "Reading, Massachusetts"
}

@book{knuth99,
    author    = "Donald E. Knuth",
    title     = "Digital Typography",
    year      = "1999",
    publisher = "The Center for the Study of Language and Information",
    series    = "CLSI Lecture Notes (78)"
}
\end{verbatim}
\end{scriptsize}
    \caption{Esimerkkejä kirjallisuuslähteiden kuvaamisesta .bib-tiedostossa.}
    \label{bibexamples-fi}
\end{figure}

Viitteet kirjallisuuslähteisiin muodostetaan komennolla \texttt{\textbackslash citep\{einstein\}}, josta generoituu tekstiin valitun viittaustyylin mukaisesti muotoiltu viite \citep{einstein}, tai \texttt{\textbackslash citep\{latexcompanion,knuth99\}}, josta tekstiin puolestaan generoituu \citep{latexcompanion,knuth99}. 
Voit esimerkiksi kirjoittaa \citep{einstein} viitataksesi julkaisuun, jonka on kirjoittanut \citeauthor{einstein} vuonna \citeyear{einstein}, kun vain lähde \citet{einstein} on oikein lisättynä kirjallisuuslähdetiedostossa (katso miltä nämä näyttävät Latex lähdekoodissa).

Tekstissä viitatut kirjallisuuslähteet tulevat automaattisesti viiteluetteloon. Kirjallisuuslähteiden tietojen oikeellisuus ja yhdenmukaisuus .bib-tiedostossa vaikuttavat luonnollisesti siihen, miten tiedot tutkielmassa näyttäytyvät. Tämä on syytä huomioida, sillä esim.\ verkosta valmiiksi {Bib\TeX} muodossa löytyvien tietojen täydellisyyten tai samanmuotoisuuteen ei pidä sokeasti luottaa.  


Keskustele viittaustyylin valinnasta ohjaajan kanssa. 
%Joitain vaihtoehtoja on osoitteessa:\\ 
%\url{https://www.overleaf.com/learn/latex/Biblatex_bibliography_styles}.
%\url{https://www.sharelatex.com/learn/Bibtex_bibliography_styles}.

\section{Joitain ohjeita Latexilla kirjoittamiseen}

Seuraavassa on joitain ohjeita tämän tutkielmapohjan käyttöön maisterintutkielmassa. Kirjoittamisohjeita löytyy useasta eri lähteestä. Voit esimerkiksi tutustua kandidaatintutkielman ohjeisiin. 
Ohjaajan kanssa on hyvä keskustella aikaisessa vaiheessa työn rakenteesta.

Yksityiskohtaisia ohjeita Latexin käyttämäsestä saa parhaiten hakemalla verkosta, esim. haku englanniksi "Overleaf table positioning" tuottaa oletettavasti aika toimivan vastauksen.

\section{Kuvat}
Kuva~\ref{fig:logo-fi} toimii esimerkkinä kuvan lisäämisestä työhön (katso tarkemmin mallia Latex lähdekoodista). Muista myös viitata jokaiseen kuvaan tekstissä. 

\begin{figure}[ht] % remove [h!] for automatic placement, which is probably better for a thesis with more text on page
\centering 
\includegraphics[width=0.3\textwidth]{template/figures/HY-logo-ml.png}
\caption{Helsingin yliopiston logo matemaattis-luonnontieteellisen tiedekunnan värein.\label{fig:logo-fi}}
\end{figure}

\newpage % just to keep the table on the same page with the short piece of text
\section{Taulukot}

Taulukossa~\ref{table:results-fi} on esimerkki kokeellisten tulosten raportoinnista taulukkona. Muista myös viitata jokaiseen taulukkoon tekstissä.
\begin{table}[ht]
\centering
\caption{Kokeelliset tulokset.\label{table:results-fi}}
\begin{tabular}{l||l c r} 
Koe & 1 & 2 & 3 \\ 
\hline \hline 
$A$ & 2.5 & 4.7 & -11 \\
$B$ & 8.0 & -3.7 & 12.6 \\
$A+B$ & 10.5 & 1.0 & 1.6 \\
\hline
%
\end{tabular}
\end{table}



% BSc instructions
%%\chapter{Johdanto}


Kaikessa julkaistavaksi tarkoitetussa tekstissä kirjoittajan luomisen ja
esitystavan vapautta rajoittavat monet ohjeet ja tarkatkin määräykset.

Parhaimmillaan lukijalle ja kirjoittajalle yhteinen, tuttu säännöstö luo
eräänlaisen tukiverkoston, joka tukee sanoman siirtymistä vääristymättä.
Kirjoituksen lukija löytää kirjoituksesta helpommin olennaisen sisällön,
jos kirjoituksen ulkoasu ja sisällön rakenne vastaavat hänen
tottumuksiaan. Sama koskee myös kirjoittajaa. Noudattaessaan valmista
esitystapamallia kirjoittajan ei tarvitse käyttää aikaansa itse työn
kannalta toissijaisten seikkojen miettimiseen, vaan hän voi keskittyä
hiomaan tekstin sisältöä. Siksi kannattaa harjoitella myös työn ulkoasua
koskevien ohjeiden noudattamista, vaikka omasta mielestään osaisikin
valita esitykselleen ohjetta paremman muodon.

Tämä kirjoitus on tarkoitettu Helsingin yliopiston
Tietojenkäsittelytieteen osastoon alempien opinnäytteiden ja
harjoitusten ulkoasun ja rakenteen ohjeeksi. Ohje soveltuu siten
kandidaatintutkielman kirjoittamisen kurssille, ohjelmistotuotantoprojekteihin, seminaareihin ja
pro gradu -tutkielmiin. (Kirjoitus on päivitetty uusintapainos aiemmista
ohjeista, jotka kurssin luennoijat ovat laatineet \citep{erkio01,erkiomakela96,erkio94,verkamo92}.)

Tyylimäärittely on saatavissa pdflatex- ja word-versiona.
Tyylimäärittelyitä valitessa on huomattava ohjeet tekstien syöttöön
liittyvässtä koodauksesta (UTF8,ISO 8859-15).
Tämän kirjoituksen tukena sopivat käytettäväksi tavanomaiset latex- tai
word-oppaat. 

\chapter{Kirjoituksen rakenne}

Tarkastellaan aluksi tieteelliseltä tekstiltä odotettuja
kirjoituksen osia. Samoihin asioihin on luonnollisesti syytä
kiinnittää huomiota myös muussa teknisessä kirjoittamisessa. Huomattakoon, että tämä teksti ei ole tieteellinen teksti, eikä siten itse sisällä kaikkia niitä elementtejä, jotka tieteellisen tekstin sisällölliseen antiin kuuluvat. Tällaisia puutteita ovat esimerkiksi johdannon tutkimuskysymyksen asettelun puuttuminen sekä arvoivan materiaalin puute tekstin lopussa, sekä yhteenvedon latteus.
Teksti rajoittuu siten otsikkonsa mukaisesti vain tekniseen sisällön asetteluun.

\section{Tiivistelmä}



Tiivistelmäsivu sisältää seuraavat osat: työn bibliografiset tiedot,
tiivistelmäteksti, aiheluokat ja avainsanat. Bibliografiset tiedot
koostuvat työn otsikosta, tekijän nimestä, julkaisupaikan tiedoista,
julkaisuajankohdasta ja sivumäärästä.

Tiivistelmäteksti on lyhyt, yleensä yhden kappaleen mittainen
(maksimissaan noin 100 sanaa) selvitys
kirjoituksen tärkeimmästä sisällöstä: mitä on tutkittu, miten on
tutkittu ja mitä tuloksia on saatu.


Aiheluokat kuvataan ACM Computing Classification System -luokituksen (CCS)
luokituksen mukaisesti. Luokittelussa käytetään täysia polkuja juurisolmun CCS osoittamista lähtöposteistä lehtisolmuihin. Polkuja voi antaa 1-3 aihepiirien soveltuvuuden mukaan, mitä alempi opinnäyte, sen vähemmän polkuja se tarvitsee. 
Poluissa tasot erotetaan toisistaan nuolella eteenpäin. Kun  polun nimisanoja arvioidaan suhteessa työn sisältöön, merkitään boldface-fontilla tärkein termi, italics-fontilla toiseksi tärkein. Näin menetellään, mikäli jotkin termeistä ovat olennaisesti paremmin kuvaavia kuin muut polun termit. Nimettyjen polkujen lisäksi lukija voi siten tarkastella lisäulottuvuutena myös tärkeiksi merkittyjen termien joukkoa sinänsä.
Avainsanoiksi valitaan kirjoituksen sisältöä
keskeisesti kuvaavia käsitteitä.

\section{Johdanto}


Johdannon tarkoituksena on kertoa yleiskielisesti
työn tavoite. Kerrotaan (kuten tiivistelmässäkin, mutta laveammin),
mitä on tutkittu, miten on tutkittu ja mitä tuloksia on saatu.
Jotta kysymyksenasettelu ja tulokset on lukijan helppo oikein tulkita on syytä aloittaa johdanto asettelemalla tutkimus asiayhteyteensä, esimerkiksi kertomalla aluksi, minkälaisessa yhteydessä tarkasteluun otettavat haasteet esiintyvät ja keiden on ratkaisuista tarkoitus hyötyä.

Johdannon pituus määräytyy suhteessa koko kirjoitelman pituuteen.
Parisivuinen kirjoitus ei erikseen otsikoitua johdantoa kaipaa, sillä
se itsessään on laajennettu tiivistelmä. Kymmensivuisen
kirjoituksen johdanto voi olla vaikkapa sivun tai puolentoista
mittainen. Pro gradu -tutkielman 50-70-sivuiseen kokonaisuuteen
tuntuu 2-4-sivuinen johdanto kohtuulliselta. 

Johdanto kertoo siis lyhyessä, yleistajuisessa muodossa
koko kirjoitelman kysymyksenasettelun, juonen sekä tulokset ja johtopäätelmät.
Tämän luettuaan lukija voi päätellä, haluaako syventyä asiaan tarkemmin
lukemalla koko kirjoituksen.


\section{Käsittelyluvut}

Käsittelylukujen työnjako määräytyy käsiteltävän asian luonteen
mukaisesti.
Lukijan ohjailemiseksi kukin pääluku kannattaa aloittaa lyhyellä
kappaleella, joka paljastaa mikä kyseisen luvun keskeisin sisältö on ja
kuinka aliluvuissa asiaa kehitellään eteenpäin.
Erityisesti kannattaa kiinnittää huomiota siihen, että lukijalle ilmaistaan selkeästi miksi kutakin asiaa käsitellään ja miten käsiteltävät asiat suhtautuvat toisiinsa. 

Jäsentelyongelmista kielivät tilanteet, joissa
alilukuja on vain yksi, tai joissa käytetään useampaa kuin
kahta tasoa (pääluku ja sen aliluvut). Kolmitasoisia
otsikointeja saatetaan tarvita joissakin teknisissä
dokumenteissa perustellusti, mutta nämä muodostavat poikkeuksen.

Perusohjeena on käyttää tekstin rakenteellisesti painokkaita paikkoja,
kuten lukujen avauksia ja teksikappaleiden aloitusvirkkeitä
juonenkuljetukseen ja informaatioaskeleiden sitomiseen toisiinsa.
Tekstikappaleiden keskiosat, samoin kuin lukujen keskiosat selostavat
asiaa vähemmän tuntevalle yksityiskohtia, kun taas aihepiirissä jo
sisällä olevat lukijat voivat alkuvirkkeitä silmäilemällä edetä
tekstissä tehokkaasti eksymättä tarinan juonesta.

Kullakin kirjoittajalla on oma temponsa, joka välittyy lukijalle
tekstikappaleiden pituudessa ja niihin sisällytettyjen ajatuskulkujen
mutkikkuudessa. Kussakin tekstikappaleessa pitäisi pitäytyä vain yhdessä
informaatioaskelessa tai olennaisessa päättelyaskelessa, muuten juonen
seuraaminen käy raskaaksi olennaisten lauseiden etsiskelyksi. Yksivirkkeisiä
tekstikappaleita on syytä varoa. 



\section{Lähdeviittausten käyttö}


Olennaisia opittavia asioita viittaustekniikoissa ovat viitteen paikka
tekstissä, oikea lähdeluettelojärjestys valitun viitetyylin parina sekä
taito ja tahto noudattaa annettua tyylimääräystä. Väitöskirjoissa ja
lehti- tai konferenssiartikkeleissa tekstin hyväksyminen riippuu myös
näiden yksityiskohtien asianmukaisesta käsittelystä. Tästä syystä
laitoksella nähdään tarpeelliseksi opiskelijoiden tutustua edes
pinnallisesti myös muihin tyylilajeihin ja oppia käyttämään
automatisoituja muotoilutyökaluja tehokkaasti, jolloin tyylimuutokset
ovat tehokkaita.

Lähdeviitteet sijoitetaan aina virkkeen sisäpuolelle. Siten esimerkiksi
tekstikappaleen lopussa irrallaan oleva viite ei ole asiallinen. Tilanne
ei muutu, vaikka viite sujautettaisiin tekstikappaleen viimeisen
virkkeen sisään. 
Lähdeviittauksen yhteyteen merkitään mukaan tarkentavat
sivunumerot, mikäli lukijan olisi työlästä löytää asianomainen kohta
viitatusta lähteestä. 


Tehokkaita viitteensijoittelupaikkoja ovat esimerkiksi uuden käsitteen
nimeämiskohta ja virkkeen loppu kun kyseessä on lähteestä lainattu
väite. On myös muistettava lainausmerkkien käyttö silloin kun tehdään
suoria lainauksia.

Tekstin jäsentelyn on tuotava selkeästi esiin, mihin asiaan viite
liittyy. Samalla tulee ymmärrettäväksi se, kuinka pitkään
tekstikatkelmaan ko. viite liitetään. Ei ole siten asiallista aloittaa
lukua nimeämällä yhtä tai useampaa lähdettä luvun taustaksi, vaan
viitteitä on kiinnitettävä täsmällisemmin väitteisiin ja käsitteisiin.
Luvun avaus viitetiedolla voi olla oire myös suuremmasta ongelmasta:
lähderiippuvuudesta. Aloitteleva kirjoittaja helposti toistaa lähteestä
oppimaansa ilman että tarpeellinen analysointi ja prosessointi suhteessa
muuhun opittuun olisi vielä tapahtut.

Viitteillä ja sanamuodolilla on
myös tuotava selkeästi esiin se, mikä teksissä on lainattua ja mikä oman
pohdinnan ja valikoinnin tulosta.


Lähdeviittauksiin käytetään Tietojenkäsittelytieteen osastolla
numeroitua tyyliä ja APA-tyyliä, valinnan näiden välillä tekevät kunkin
ryhmän valvoja ja ohjaaja yhdessä. 
Numeroitu tyyli on esimerkiksi IEEE- ja
ACM-julkaisuissa yleisesti käytetty ja puolustaa siten paikkaansa.
APA-tyyli on poikkeuksellinen ns. kovissa tieteissä, mutta monet
valvojista pitävät siitä sen luettavuuden vuoksi. Numeroita joutuu nimiä
useammin tarkistamaan lähdeluettelosta, sillä tarkastus- ja
arvointiprosessiin kuuluu arvioida myös lähteiden valitaa ja niiden
käyttötapaa.





\section{Yhteenveto}

Yhteenveto  vaatimattomimmillaan on vain lyhyt kertaus kirjoituksen
keskeisistä asioista. Arvokkaamman yhteenvedon saa aikaan kommentoimalla
 työn tulosten arvoa, työn liittymistä ympäristöön ja
tulevaisuudennäkymiä. Tällaiset arviot  huolellisesti
perusteltava.

\section{Lähdeluettelon laatiminen}

Tieteellisen kirjoittamisen kurssin töiden lähdeluetteloiden
laatimisessa noudatetaan seuraavia ohjeita.

Niiden taustalla on kaksi
keskeistä pyrkimystä: tehdä viitatun lähteen hankkiminen luettavaksi
mahdollisimman helpoksi ja ilmaista, millaisen arviointiprosessin
läpi käyneeseen kirjoitukseen vedotaan.
Näistä syistä
\begin{itemize}
\item lähdeviitteen tulee aina olla niin tarkka, että
lähde on sen perusteella tunnistettavissa ja löydettävissä luetteloista
ja kirjastoista,
\item erityyppisten lähteiden (kirjat, konferenssit, lehdet) on erotuttava
toisistaan
ja
\item luettelon eri osien tulee olla mahdollisimman
yhdenmukaisia, erityisesti lähdetyypin sisällä.
\end{itemize}


Riippumatta käytettävästä viitetyylistä, 
lähteet ovat Tietojenkäsittelytieteen osaston opinnäytteiden lähdeluetteloissa tekijän nimen mukaisessa aakkosjärjestyksessä,
saman tekijän (tekijäryhmän) työt julkaisuajan mukaisessa
järjestyksessä. Jos jollakin lähteellä ei ole henkilötekijää, se
aakkostetaan julkaisun nimen mukaisesti. 

Kustakin lähteestä annetaan seuraavat tiedot, edelleen viitetyylistä riippumatta:
\begin{itemize}
\item (tarvittaessa lähdeviitelyhenne).
\item
tekijän tai tekijöiden nimet (sukunimi, etunimien alkukirjaimet)
alkuperäisessä järjestyksessään; jos tekijöitä on enemmän kuin kolme,
voidaan toimia siten, että
vain ensimmäinen tekijä nimetään ja muiden tilalle kirjoitetaan {\em et
al.}
\item
julkaisun tai artikkelin nimi alkuperäisessä muodossaan
\item
julkaisupaikan tiedot:
\begin{itemize}
\item
kirjasta: kustantaja, julkaisupaikka (voidaan jättää pois, jos kyseessä
on tunnettu kustantaja), vuosi ja
\item
lehtiartikkelista: lehden nimi, volyymi, numero, vuosiluku ja kuukausi (suluissa),

\item
artikkelikokoelmassa (esim. konferenssijulkaisussa) ilmestyneestä
artikkelista:
\begin{itemize}
\item kokoelman nimi, toimittaja, kustantaja, julkaisupaikka ja vuosi
{\em tai}
\item konferenssin nimi, järjestäjä, paikka ja aika,
\end{itemize}
\item
raportista: julkaisusarja, raportin numero, julkaisupaikka, julkaisija ja vuosi
ja
\item
www-lähteestä: verkko-osoite, voimassaoloajankohta, mahdollisesti
viittausajankohta hakasuluissa
\end{itemize}
\item
sivunumerot, mikäli lähteenä käytetty julkaisu on artikkeli tai kokoomateoksen itsenäinen luku.
\end{itemize}

Normaaliin suomalaiseen tapaan artikkelin nimessä ainoastaan
ensimmäinen sana kirjoitetaan isolla alkukirjaimella, sen sijaan
konferenssien ja kokoelmajulkaisujen nimissä käytetään isoa
alkukirjainta jokaisen sanan alussa (artikkelisanoja ja prepositioita
lukuunottamatta). Katso mallia oheisista esimerkeistä.
Kokoelman nimen edessä on syytä selvyyden vuoksi käyttää sanaa {\em
Teoksessa}, paitsi kun on kysymys konferenssijulkaisusta, jonka nimi
alkaa lyhenteellä {\em Proc.} (sanasta Proceedings). Tällöin ei tarvita
mitään täydennystä.
Tämän eron näkee esimerkiksi vertaamalla
lähdeviitteiden~''\citep{dantowsley90}''
ja~''\citep{gannonetal89}'' ulkoasuja.

WWW-lähteiden käytössä on syytä muistaa, että verkossa julkaisukynnys on
olematon. Kannattaa siten keskittyä tunnettujen tieteellisten
kustantajien julkaisuihin ja niihin teknisiin standardeihin, joille WWW
on ainoa julkaisukanava. Mikäli sama julkaisu on saatavissa myös
perinteisessä muodossa, viitataan ensisijaisesti siihen ja käytetään
verkko-osoitetta lisätietona. Lähdeluettelossa on annettu esimerkit
useita kanavia julkaistusta kirjoituksesta~\citep{abiteboul,dietinger} sekä pelkästään
WWW-julkaisuna
leviävästä standardista~\citep{bray}.

Erityisesti varoitetaan Wikipedian käytöstä tieteellisessä tekstissä.
Vaikka sen avulla on helppo alustavasti tutustua joihin aihepiireihin ja
asiantuteva lukija voisi teksin kelvolliseksi tiettynä hetkenä
hyväksyäkin, ei se foorumina millään lailla täytä tieteellisesti
vertaisarvoidun tutkimusfoorumin kriteerejä.
Jos Wikipedia-artikkelia  ei mitenkään malta ajankuvana olla mainitsematta, käytettäköön jotain muuta kuin lähdeviitetekniikkaa tähän taiteelliseen otteeseen, vaikkapa alaviitteitä. Olennaista silloinkin on, että tieteellinen sisältö ei tule tällä korvatuksi vaan sen puute korostetuksi.

WWW-lähteeseen viittaamisessa pätevät samat periaatteet kuin
perinteisiin lähteisiin viitattaessa: lähdeviitteessä ilmaistaan
otsakkeet, kirjoittajat, toimittajat js muut seikat. Eroa on ainoastaan
verkko-osoitteen ja sen voimassaoloajankohdan ilmaisemisessa. Mikäli
lähde on julkaistu ainoastaan verkossa, voidaan web-osoitetta (URL)
käyttää vastaavasti kuin perinteisen julkaisun paikannusinformaatiota
(lehden ja se numeron julkaisutiedot). Lähdeluettelossa on WWW-viittausten yhteydessä aina syytä ilmaista päivämäärä, jolloin linkin voimassaolo ja lähteen sisältö on tarkastettu.
Esimerkkeinä verkkoviitteistä soveltuvat seuraavat:
\begin{itemize}
\item Gergen, Kenneth (1999) Narrative, Moral Identity and Historical
Consciousness: a Social Constructionist Account.
http://www.swarthmore.edu/SocSci/kgergen1/text3.html. Haettu 11.6.1999.
\item 	
Ritala-Koskinen, Aino and Valokivi, Heli (2006) The Role of Development
Skills in Social Work Practice Education in Finland. Social Work and
Society, The International Online-Only Journal 4(2006)1.
http://www.socwork.net/2006/1/series/transition/ritalakoskinenvalokivi.
Viitattu 30.8.2006.
\item Heinisuo, Rami and Ekholm, Kai (1997) Elektronisen viittaamisen
opas. Jyväskylän yliopiston kirjaston julkaisuja n:o 40. Jyväskylä:
Jyväskylän yliopiston kirjasto. http://www.pori.tut.fi/~multisil/evo/.
Viitattu 29.8.2006.
\end{itemize}

Kirjoituksen lähdeluettelossa luetellaan täsmälleen ne lähteet, joihin
viitataan kirjoituksen tekstiosassa. Tämän kirjoituksen lähdeluettelo on
tarkoitettu lähinnä esitystavan esimerkiksi, mistä syystä siinä on
''ylimääräisiä'' lähteitä.


Pääsääntöisesti julkaisun tai artikkelin nimen perään tulee piste,
samoin kunkin lähteen bibliografisten tietojen perään. Muut tiedot
erotetaan toisistaan pilkulla. Useimmissa tapauksissa  
voidaan noudattaa teknisten välineiden antamaa mallia, sillä edellytyksellä, että ylläolevat vaatimukset muuten täyttyvät.



\chapter{Ulkoasulliset seikat}

  Tässä luvussa käsitellään yleisimpiä
tekstin tekniseen esittämiseen liittyviä seikkoja.  

\section{Työn osien järjestys}

   Kirjoituksen alussa on aina
erillinen, mallin mukainen kansilehti. Toisena sivuna on
tiivistelmäsivu, sen jälkeen sisällysluettelo (yksi tai useampia
sivuja) ja sitten varsinainen teksti. Sivunumerointi aloitetaan vasta
ensimmäiseltä tekstisivulta (arabialaisella ykkösellä). (Tarkat
jättävät ykkössivun numeromerkittä.) Sisällysluetteloon merkitään
kaikki (numeroidut) otsikot ja vastaavat sivunumerot. Monet
tekstin\-käsittelyjärjestelmät muodostavat itse sisällysluettelon,
jolloin kirjoittajan ei tarvitse huolehtia luettelon sivunumeroiden
päivittämisestä tekstin kehittyessä. Sisällysluettelo\-sivu ja sitä
edeltävät sivut voidaan haluttaessa numeroida erikseen (roomalaisin
numeroin) esimerkiksi tämän mallin mukaisesti.  

Varsinaisen tekstin
jäljessä, mutta itse työhön kuuluvana on ensimmäisenä lähdeluettelo,
jonka otsikkoa ei numeroida. Lähdeluettelon jälkeen sijoitetaan
mahdolliset liitteet, jotka otsikoidaan ja varustetaan sisäisillä
sivunumeroilla.  



Mikäli kuvista, algoritmeista
ja taulukoista halutaan tehdä yhtenäinen luettelo, sijoitetaan
luettelot sisällysluettelon jälkeen. Luetteloiden käyttöarvosta on
eriäviä mielipiteitä, joten niiden laatimiseen ei varsinkaan ilman
tekstinkäsittelyjärjestelmän tukea kannata ryhtyä ilman tarkastajan
erityistä toivetta.  

Mikäli kirjoitukseen erityissyistä halutaan
liittää aakkosellinen hakemisto, sijoitetaan se lähdeluettelon jälkeen
ennen liitteitä. Indeksi merkitään sisällysluetteloon samoin kuin
lähdeluettelo (numeroimaton luku). Mikäli indeksin tekemiseen
ryhdytään, on syytä käyttää tekstinkäsittelyjärjestelmän tarjoamaa
automatiikkaa.

Teksin luonnollisen juonenkuljetuksen mukana esiin
tulevien käsitteiden määrittelyjen sijasta ei pidä yrittää sen enempää
pakata kaikkia määritelmiä johdantoon kuin laatia johdantoa ennen
käsitelistaa tai lyhenteiden selityslistaa. Kumpikaan ei sovi
tavanomaiseen argumentoivaan tieteelliseen tekstityyliin, vaikka
teknisessä yhteydessä niillä liitteinä voi olla lisäarvoa.


\section{Tekstin yleinen sijoittelu}

Lopullinen tutkielmaversio voi olla yksi- tai kaksipuoleiseksi aseteltua
ja riviväliltään 1,5 tai 1.  Erityyppisissä
teksteissä haasteet ja asetteluvaatimukset voivat olla erilaiset. Erota
kappaleet toisistaan yhdellä tyhjällä rivillä tai
käytä tekstinkäsittelytyökalujen ominaisuuksia
hyödyksesi ja määrittele tekstikappaleiden väliin jäävä tila hieman
normaalia riviväliä suuremmaksi.

Kirjoituksen lukujen, kuvien ja taulukoiden erottumisen kannalta
tärkein keino on riittävän tilan käyttö niiden ympärillä. Kuvan ja
nimekkeen tulee olla selkeästi yksi kokonaisuus, joka eroaa muusta
tyhjän tilan rajaamana. Kuvan tai taulukon on aina numerointinsa ja
nimekkeensä kanssa mahduttava yhdelle sivulle tai varmasti
kaksipuolisena paperidokumenttina tarkasteltavassa tekstissä aukeamalle.
Kuvissa fonttikoko ei saa alittaa 8 pistettä.



Jos uusi luku tulisi alkamaan aivan sivun alareunasta (vain yksi tai
kaksi riviä varsinaista tekstiä), aloita mieluummin uusi sivu. Jokaista
uutta lukua ei kuitenkaan ole tarpeen --- etenkään lyhyessä
kirjoituksessa --- aloittaa uudelta sivulta: jos kirjoituksessa on
paljon melkein tyhjiä sivuja, lukija voi epäillä, että kirjoittaja on
yrittänyt saada kirjoituksensa näyttämään pitemmältä kuin se onkaan. 
 Tyhjää tilaa kannattaa käyttää hyödyksi myös kuvien ja taulukoiden
yhteydessä. Erityisesti jos kirjoituksessa käytetään kauttaaltaan samaa
tekstityyppiä, tyhjät rivit ovat välttämättömiä erottamaan esimerkiksi
tekstiä ja taulukkoa toisistaan. Tyhjä tila on halpaa, mutta se lisää
selkeyttä ja luettavuutta.  


\section{Kuvat ja taulukot}


Kuva tai taulukko sijoitetaan mahdollisimman lähelle
(ensimmäistä) tekstikohtaa, jossa siihen viitataan, ei kuitenkaan
kyseistä viittausta aikaisemmaksi.
Tekstissä on syytä myös kertoa, mitä kuvalla halutaan havainnollistaa.
Kuvan voi lukea monella eri tavalla, joten lukijaa on ohjattava.

Kuvaa ei pidä sijoittaa välittömästi luvun otsikon alle, vaan on
aloitettava tekstillä. Kuvaa ei pidä sijoittaa keskelle tekstikappaletta
(saati virkettä), paitsi jos kuva tulee sivun alkuun tai loppuun eikä
kappaleen jatkumisesta tule epäselvyyttä.

Kuvan ei aina tarvitse olla välittömästi viittaavan kappaleen
perässä. Esimerkiksi viittauskohdan ja
vasta seuraavalle sivulle mahtuvan kuvan väliin jäävää sivun loppuosaa
ei jätetä tyhjäksi. Kuvaa ei kuitenkaan pidä viedä seuraavaa
sivua kauemmas viittauskohdasta.


Varsinaista kuvan esittämistä havainnollistaa kuva~\ref{kuvaesimerkki}.
Huomiota on kiinnitettävä kuvan osien ja tekstimerkintöjen näkyvyyteen,
kuvan numerointiin ja otsikointiin. 

\begin{figure}[ht]
%\begin{figure}[tbh] t= top, b = bottom, h=here
\ \newline
\begin{center}
\includegraphics[width=0.75\textwidth]{kuvaesimerkki.pdf}
\caption{Kuvan elementit.}
\label{kuvaesimerkki}
\end{center}
\end{figure}


Kuvien kokoon on kiinnitettävä huomiota. Käytettyjen merkintöjen
on oltava helposti luettavissa ja selkeät. Esimerkiksi
suorituskykykäyriä esitettäessä akselit on nimettävä, asteikot
merkittävä ja käytetyt yksiköt tuotava selkeästi esiin.
Samankaltaisia asioita esitettäessä useammalla kuvalla on
syytä käyttää samaa mittakaavaa vertailun helpottamiseksi.

Kuvan otsikko kirjoitetaan kuvan alle ja sen tulee olla mieluummin lyhyt
ja ytimekäs kuin liian selittelevä.
Samoin toimitaan taulukoiden otsikoinnissa.

Kuvat ja taulukot numeroidaan juoksevasti. Pitkissä teksteissä käytetään
kaksitasoista numerointia (esimerkiksi Kuva 3.1) pääluvuittain, lyhyissä
riittää yksitasoinen numerointi.


Kuva- ja taulukko-otsikoiden yhdenmukaiseen esitystyyliin on syytä kiinnittää
huomiota, samoin mm. välimerkkeihin. Luontevaa on käyttää
kuvatekstin lopussa pistettä, ovathan useimmat kuvateksteistä virkkeitä. 

(Kuvien ja taulukoiden otsikointityyli vaihtelee
kustantajittain ja julkaisuittain. Samoin tuntuu suositeltava käytäntö
Tietojenkäsittelytieteen laitoksen sisällä vaihtelevan taulukon
otsikon sijainnin suhteen.)


\section{Otsikot}

Otsikoissa voi käyttää muusta tekstistä poikkeavaa kirjasintyyppiä,
alleviivausta, suurempaa kirjasinkokoa tms.\ erotuskeinoa, yleensä
kuitenkin vain yhtä näistä, koska kovin monta erilaista kirjasintyyppiä
ja -kokoa tekee ulkoasusta helposti sekavan.  Otsikoiden esitystavan on
oltava johdonmukainen läpi koko kirjoituksen. Numeroimattomia
''ylimääräisiä'' otsikoita ei tule yleensä käyttää.


\section{Mallin käyttö}

Voit käyttää tätä kirjoitusta mallina oman opinnäytteesi ulkoasua
varten. Eri tekstinkäsittelyjärjestelmissä käytössä olevat yksityiskohdat kuten
kirjasintyypit ja -koot ja rivivälit  poikkeavat toisistaan, joten
pienet poikkeamat ovat toki hyväksyttäviä.

Tieteellisen kirjoittamisen kurssin luennoilla ja
liitteenä olevassa ohjeessa annetut töiden ohjeelliset sivumäärät
koskevat työtä, joka vastaa ulkoasultaan tätä ohjetta (kirjasinkoko
12~pistettä). Tässä tekstissä keskimääräinen rivin pituus lienee noin
80~merkkiä ja sivun pituus 35-40~riviä.
Sivumääriin lasketaan varsinaisen tekstiosuuden pituus ja lähdeluettelo
(arabialaisin numeroin numeroitu osuus), ei kansilehteä, tiivistelmää
eikä sisällysluetteloa. Sivumääräarviossa otetaan huomioon hyvin vajaat
sivut, joita syntyy paljon lyhyiden lukujen ja taittotyyliin määritellyn
luvun avauksen pakottaminen oikeanpuolimmaiselle sivulle. 

\chapter{Yhteenveto}

Tämän kirjoituksen tarkoituksena on toimia muistilistana eräistä
esitystavallisista säännöistä, joihin harjoitusten ja tutkielmien
kohdalla on syytä kiinnittää huomiota.

Annetut ohjeet on laitoksen henkilökunta muotoillut yhdessä keskustellen
ja noudattaen oman tieteenalansa perinteitä. Eri erikoistumisaloilla ja
erilaisillaa määräävässä asemassa olevissa julkaisufooruilla käytänteet
vaihtelevat ja nuorten tutkijoiden onkin tiedostettava ero yleisten
sisältöohjeiden ja teknisten muotoilusääntöjen välillä. Aina tekstin
valmistuessa on tarkastettava erikseen, täyttääkö se annetut
pituusrajoitteet ja vastaako se annettuja muotoiluohjeita, olivatpa ne
kuinka pikkutarkkoja tahansa. Tarkasta sääntöjen noudattamisesta syntyy
yhteinäisyyttä kokoovan julkaisun tasolla, mikä helpottaa lukijoiden
työskentelyä.

Tämä ohje vastaa vain asettelullisiin kysymyksiin ja sen rinnalla on
syytä tutustua materiaaliin ja luentoihin, joissa keskitytään tekstin
varsinaiseen sisältöön. Olennaisin väline on kuitenkin akateemisesti
pidemmälle ehtineen, jo julkaisuja rakentaneen ohjaajan palaute ja
mentorointi.


%%
\chapter{Introduction}


In all writing for publication, the writer's freedom of creation and expression are limited by a number of guidelines and specific regulations.

At best, a familiar set of regulations shared by reader and writer can create a kind of support network that allows the message to be relayed without distortion. It will be easier for readers to find 
the pertinent contents in a piece of writing if its layout and structure are the same as they are used to. This also applies to writers. When writers follow a set presentation model, 
they do not have to waste time on considerations that are secondary to the work itself, but they can concentrate on polishing the contents of the text. This means that it is a good 
idea to practice following the rules for the layout, though you may think you know how to select a better way to present your work.

This is a guide for the layout and structure of theses and essays at the Department of Computer Science at the University of Helsinki. It is thus applicable to the course 
Scientific Writing, the software engineering projects, seminars, and MSc theses. (This is an updated version of 
the previous guide written by the course lecturers \citep{erkio01, erkiomakela96, erkio94, verkamo92}.)

The \LaTeX\ guide and \LaTeX\ style that has been  published on the department's web site can be used as support for this guide.


\chapter{Structure}

Let us start by looking at the sections expected to be in a scientific text. Keep in mind that the same expectations go for all kinds of technical writing.
However, this document in itself, is not a scientific or a research
text, so there will be content lacks in terms of research question
setting in the introduction and evaluative material in the last sections.

\section{Abstract or summary}
%\enlargethispage{5mm}


The summary page contains the following elements: the bibliographical data of the work, an abstract, topic classification, 
and the key words. The bibliographical data consists of title, name of the author, place of publication, date of publication, and number of pages.

The abstract should be short, generally one paragraph (100 words maximum) explaining the main contents of the work: topic, methodology and results.


 Topics are classified according to the ACM Computing Classification System
(CCS). A small set of paths (1-3) should be used, starting from any top nodes
referred to bu the root term CCS leading to the leaf nodes. The elements
in the path are separated by right arrow, and emphasis of each element individually can be indicated
by the use of bold face for high importance or italics for intermediate
level. The combination of individual boldface terms may give the reader
additional insight. 

\section{Introduction}


The purpose of the introduction is to introduce the goals of the work in general terms. Describe topic, 
methodology and results (as in the abstract, but expand it).
In order to provide the reader a good starting point for
interpretations, it is good to start the introduction by
contextualisation of the challenges and solutions to be discussed. For
example, why a certain domain has a particular challenge and who are
intended to benefit from the solutions proposed.

The length of the introduction depends on the length of the whole work. A few pages of text does not need a 
separate introduction, since it is an expanded summary in itself. The introduction to a 10-page text can 
be 1--1.5 pages long. For a 50--70-page MSc thesis, a 2--4-page introduction seems reasonable.

The introduction should shortly describe the problem field of the whole work, the plot, and the conclusions, 
in general terms. After reading it, the reader may decide whether to go deeper into the topic by reading the whole text.


\section{Topic chapters}

The nature of the matter at hand determines how the topic chapters are disposed.
In order to guide the reader, it is a good idea to start each main chapter with a short paragraph 
on what the main topic of the chapter is and how it progresses from one sub-chapter to the
next. Especially relevant is to express how concepts, challenges,
solutions and research steps are bound together. There should be enough
guidance for the reader to allow expectation of the right storyline.


Basic rule for easy to follow text is to use the natural emphasis of the
text structure to support the content matter key concepts and thought
processes. This means using openings of sectiosn and text paragraphs for
key arguments and information moves, while the internal parts of
paragraps are filled with supporting aspects that those less familiar
with the topic area need. Those with better background knoweldge can
quickly skim trhough the text without loosing any essential arguments.

Each author has his or her personal ryhtm in the text, which is visible
for the reader as the length of text paragraphs and complexity of
thought chains within. A good policy is to take only one information
move or transition per paragraph. This way the text stays easy to follow.

Signs of problems with the disposition of the text are easily seen in texts
with only one sub-chapter, or with more than two chapter levels (main and sub-chapters). There may be justifiable reasons to use three-level 
headings in some technical documents. Single sentence text paragraphs
are also to be avoided.



%\pagebreak
\section{Reference usage}


Relevant learning targets include superficial knowledge of several
citation styles and capability (and willingness) to follow a given
style and ordering of entires and bibliographical details in the list of references.
These aspects are essential as the approval of a PhD manuscript or
journal article may depend on them.

Disregard of which style you use, 
references are always placed inside sentences.  This means that e.g. a separate reference at the end of a paragraph would be inappropriate.

The structure of the text must clearly show what the reference relates to.  At the same time, it 
shows how long a piece of the text that the reference relates to.

Efficient positions for citations are right after the introduction
(definition) of a concept, a methods or such, or in the end of a claim
from the reference material. Furthermore, if quoting verbatime one must
use citation marks.

The text structures and wordings, in addition to the location of
ciations must clearly express whether claims or arguments are 
authors' own, or if they come from the contribution of others. Thus is
is of bad style to open a section by listing the citations on which the
section is based on. Such method can further indicate more serious
problems like following reference material as it is instead of analysing
and synthetising material into new though processes. 


\section{Conclusion}

At its simplest, a conclusion is merely a weak revision of the main points of the text.  All more valuable 
conclusion sections contain comments on e.g. the value of results, how the work relates to its environment, or 
future visions. This kind of evaluations should be well-grounded in fact, though, or the conclusion 
might inadvertently seem comical. 

\section{Creating the list of references}

The following guidelines should be followed when creating lists of reference for the 
assignments during the course Scientific Writing.

The guidelines are backed by two main goals: to make it as easy as possible to find the 
referenced source, and to show what kind of evaluation process the referenced work has undergone. For these reasons
\begin{itemize}

\item the reference notes should always be so exact that the source can be recognized and found in catalogues and libraries

\item different types of sources (monographs, conferences, journals) have to be easy to distinguish from each other, and 

\item the different parts of the list must conform to each other, especially for each source type.
\end{itemize}


Independent of the citation style in use, 
the sources are listed alphabetically according to the author's name, and works by the same 
author (group of authors) in the order in which they have been published. If some publication 
does not have an individual writer, it is alphabetized according to its name.


The following information should be given on each source independent of
the citation style:
\begin{itemize}
\item
Name(s) of author(s) (surname, initial letters of first names) in the original order; if there are more 
than three authors, we can write the name of the first author and {\em et al.} instead of the other names.

\item
the name of the publication or article in its original form

\item
place of publication:

\begin{itemize}

\item
of monographs: publisher, place of publication (can be omitted if the publisher is well known), year
\item
of journal articles: name of journal, volume, issue, year and month (in parenthesis)

\item
of articles from article collections (such as conference publications):
\begin{itemize}

\item name of collection, editor, publisher, place of publication and year 

{\em or}

\item conference name, coordinator, place and time,

\end{itemize}

\item
of a report: series, report number, place, publisher and year

\item
of a web source: URL, validity date, possibly the date when referenced in square brackets

\end{itemize}

\item
page numbers, if the source is an article or constitutes a chapter in a compilation.
\end{itemize}


When using web sources, you should keep in mind that the threshold for publication on the web 
is non-existent. It is better to concentrate on the publications of well-known scientific 
publishers and the technical standards for which the web is the only publication channel. If 
the same publication is available on paper, refer to that primarily and add the URL as a complement.


The list of references gives an example of a text that has been published through many 
channels~\citep{abiteboul,dietinger}and another example that shows a standard that has 
been disseminated only through the web~\citep{bray}. For web-based
references it is important to give both the publication date and the
time of reading and interpreting the material. The content is prone to
change, and dating your use, you protect your interpretations for
unnecessary accusations of being faulty, if it is known which version
you used.



The list of references for a text should list exactly the sources that the text refers to. 
The list of references for this text is an example of how to present sources, 
which is why it contains ''extra'' sources.


In any style,
put a full stop after the name of a publication or article, as well as after the bibliographical 
data of each reference. Separate the other pieces of information with a comma. As is normal in 
Finnish, only the first letter of the first word in the heading is capitalized, but in the titles of 
conferences and compilations, each major word is capitalized (not articles or prepositions). See the 
appended example for a model. For the sake of clarity, it is best to write {\em In the work} before 
the name of a compilation, except in the case of conference publications where the name starts with 
the abbreviation {\em Proc.} (for Proceedings). In such cases, no complement is necessary. You can 
see the difference by comparing the layouts of the references ~''\citep{dantowsley90}'' and~''\citep{gannonetal89}'' .
In case of other citation styles, it is likely that you can trust the
results given by the automated reference management tools.

\chapter{Layout}

This chapter discusses the main issues of the technical presentation of a text.


\section{Disposition of text sections}


Each text should include a separate cover page, as in this model. The second page contains an 
abstract, followed by a table of contents (on one or more pages), and then the main text.  
Pagination starts on the first page of the main text (with the Arabic numeral 1). (Rigorous writers 
leave out the number from the first page.) All the (numbered) headings and their page numbers should be 
written into the table of contents. Many word-processing systems create the table of contents automatically 
so that the writer does not have to worry about updating page numbers as the work progresses. If writers 
wish to, they can paginate the table of contents and previous pages separately (with Roman numerals), e.g. like this model.

After the main text, but as part of the text body, comes the list of references; its heading is 
not numbered. Any possible appendices are added after the list of references, with headings and internal page numbers.

%\enlargethispage{5mm}

If you want to make a coherent list of figures, algorithms and tables, this list should be placed i
mmediately after the table of contents. The value of such lists is a matter of opinion, so there 
is no need to create one --- especially if your word-processing software does not support it --- unless your instructor explicitly asks for it.

If for some special reason you want to add an alphabetical index top your work, it should be placed after 
the list of references and before appendices. An index should be noted in the table of contents, as s
hould the list of references (unnumbered chapter). Writers who want to create an index should use the automatic support their word-processing system offes.

\section{General text layout}

Nowadays you may plan your text to be printed on single side or double
side, with line spacing one, for the final copy.

For reviewing and feedback purposes, check with your readers for their
preferences. Most read and comment on paper and need some space for
marking. 
For drafts you can use wide margins too. 
Both margins should also be fairly broad (ca.~3~cm): the wider left margin is needed for binding the text and the
right one for the instructor's comments.  Leave enough space (2--3~cm) at the top and bottom, as well.

The most effective way to distinguish features like chapters, figures, etc,  is to have enough space 
in your text.  Separate paragraphs from each other with one and chapters with a few empty lines.  If a 
new chapter is about to start at the bottom of a page (with only one or two lines of text), it is better 
to start it on the next page.  However, it is not necessary to start every chapter on a new page, especially 
not with a short text; if the text contains many pages that are nearly empty, readers might suspect that the 
writer has tried to make it look longer than it is.


Any empty space may be utilized for displaying figures and tables.  Especially if a text is written with the 
same type of text throughout, empty lines are necessary for separating e.g.~text from tables.  Empty space 
is cheap, but adds to the clarity and readability of a text.

\section{Figures and tables}

%\enlargethispage{5mm}

All figures and tables should be placed as near the (first) place in the text that refers to them, but 
not before this reference. The text should also explain what the writer wants to illustrate with the 
figure. Figures can be interpreted in many different ways, so the reader needs guiding.

Figures should never be placed immediately under the heading of a chapter, but chapters should always 
start with text. Figures should not be placed in the middle of a paragraph (much less a sentence), except 
if the figure is placed at the top or bottom of a page and it is clear that the paragraph continues.

Figures do not always have to be placed immediately after the paragraph that refers to them. If there is 
not enough space for a figure on the same page as the paragraph that refers to it, the rest of the page 
should not remain empty, Though the figure is inserted on the following page. However, the figure should 
never be more than one page away from the reference.


The image~\ref{kuvaesimerkki} shows how to present a figure. You must pay attention to the visibility of 
figure parts and text, to the numbering of figures, and captions. 

\begin{figure}[ht]
%\begin{figure}[tbh] t= top, b = bottom, h=here
\ \newline
\begin{center}
\includegraphics[width=0.9\textwidth]{kuvaesimerkki.pdf}
%\rotatebox{90}{\includegraphics[scale=.75]{kuvaesimerkki.pdf}}
\caption{Figure elements.}
\label{kuvaesimerkki}
\end{center}
\end{figure}


We must also pay attention to the size of figures. Any annotations must be clear and easy to read When 
presenting performance curves, for example, the axels have to be named, the scale notated, and the units 
clearly presented. If you present many things with similar figures, you should use the same scale for easy comparison.

The caption of a figure should be written underneath it, and it is preferable that it be short and to 
the point than too explicit. The same goes for table captions.

Figures and tables should be numbered progressively. In long texts, two-level numbering should be 
used (e.g. Figure 3.1) according to the chapter number, but in shorter texts, one-level numbering is good enough.


You should pay attention to presenting figure and table captions consistently, as well as to punctuation 
marks. It is natural to put a full stop after a caption, since they are most often full sentences. 

(The style for figure and table captions vary according to publisher and publication. The recommendations 
at the Department of Computer Science also seems to vary with respect to where the caption should be placed.)


\section{Headings}

You can use a different font in headings than in the rest of the text, or underlining, larger fonts, or other 
methods of emphasis, but generally it is preferred that you only use one of these methods, because if there are 
very many font types and sizes, the layout looks messy.  The format of headings must be consistent throughout 
the text. You should not use any unnumbered ''extra'' headings.


\section{Using this model}

You can use this text as a model for the layout of your own assignment. The font types and sizes, 
line spacing, etc vary according to word-processing system, so small deviations from the rule are acceptable.

The directive number of pages for written assignments given during the lectures for the Scientific-writing 
course and in this guide are applicable to assignments with a layout like this guide's (font size 12 points). The 
average line in this text should consist of about 80~characters and one page about 30~lines. The number of pages 
includes the text itself and the list of references (the part paginated with Arabic numbers), not the 
cover page, summary, or table of contents.

\chapter{Conclusion}

This text is a checklist for some of the rules governing written presentations, which you should 
keep in mind when writing exercises and theses.

This collection of advise has been compiled by staff members as the
result of discussions on general good writing habbits in computer
science. The consensus is that all young researchers and academic degree
holders must be able to seek and follow instructions of layout and
structuring of texts depending on the contribution they are working on. 
This set of instructions aims at unifying the looks of theses and other
study reports from the department, but it is also a representative of
the instruction sets students will meet later.

This set of instructions only address the looks of the document. Other
instructions must be sought, read and gained experience with in order to
learn to select the suitable semantical contents for the scientific texts.



\end{appendices}
%%%%%%%%%%%%%%%%%%%%%%%%%%%%%%%%%%%%%%%%%%%%%%%%%%%%%%%%%

\end{document}
